\PassOptionsToPackage{unicode}{hyperref}
\PassOptionsToPackage{hyphens}{url}
\PassOptionsToPackage{dvipsnames,svgnames,x11names}{xcolor}
\documentclass[10pt,letterpaper]{article}

\usepackage{xcolor}
\usepackage[margin=2.0cm]{geometry}
\usepackage{lmodern}
\usepackage{amsmath,amssymb,amsfonts,mathtools}
\usepackage{textcomp}
\usepackage{microtype}
\usepackage{parskip}
\usepackage{longtable,booktabs,array,calc}
\usepackage{caption}
\captionsetup[table]{skip=6pt,font=small,labelfont=bf}
\captionsetup[figure]{font=small,labelfont=bf}
\newcounter{none}
\usepackage{etoolbox}
\usepackage{graphicx}
\usepackage{fancyvrb}
\newcommand{\VerbBar}{|}
\newcommand{\VERB}{\Verb[commandchars=\\\{\}]}
\DefineVerbatimEnvironment{Highlighting}{Verbatim}{commandchars=\\\{\},fontsize=\footnotesize}
\RecustomVerbatimEnvironment{verbatim}{Verbatim}{fontsize=\footnotesize}
\newenvironment{Shaded}{}{}
\newcommand{\AlertTok}[1]{\textcolor[rgb]{0.80,0.00,0.00}{\textbf{#1}}}
\newcommand{\AnnotationTok}[1]{\textcolor[rgb]{0.38,0.63,0.69}{\textbf{\textit{#1}}}}
\newcommand{\AttributeTok}[1]{\textcolor[rgb]{0.49,0.56,0.16}{#1}}
\newcommand{\BaseNTok}[1]{\textcolor[rgb]{0.25,0.63,0.44}{#1}}
\newcommand{\BuiltInTok}[1]{\textcolor[rgb]{0.00,0.50,0.00}{#1}}
\newcommand{\CharTok}[1]{\textcolor[rgb]{0.25,0.44,0.63}{#1}}
\newcommand{\CommentTok}[1]{\textcolor[rgb]{0.38,0.63,0.69}{\textit{#1}}}
\newcommand{\CommentVarTok}[1]{\textcolor[rgb]{0.38,0.63,0.69}{\textbf{\textit{#1}}}}
\newcommand{\ConstantTok}[1]{\textcolor[rgb]{0.53,0.00,0.00}{#1}}
\newcommand{\ControlFlowTok}[1]{\textcolor[rgb]{0.00,0.44,0.13}{\textbf{#1}}}
\newcommand{\DataTypeTok}[1]{\textcolor[rgb]{0.56,0.13,0.00}{#1}}
\newcommand{\DecValTok}[1]{\textcolor[rgb]{0.25,0.63,0.44}{#1}}
\newcommand{\DocumentationTok}[1]{\textcolor[rgb]{0.73,0.13,0.13}{\textit{#1}}}
\newcommand{\ErrorTok}[1]{\textcolor[rgb]{0.80,0.00,0.00}{\textbf{#1}}}
\newcommand{\ExtensionTok}[1]{#1}
\newcommand{\FloatTok}[1]{\textcolor[rgb]{0.25,0.63,0.44}{#1}}
\newcommand{\FunctionTok}[1]{\textcolor[rgb]{0.02,0.16,0.49}{#1}}
\newcommand{\ImportTok}[1]{\textcolor[rgb]{0.00,0.50,0.00}{\textbf{#1}}}
\newcommand{\InformationTok}[1]{\textcolor[rgb]{0.38,0.63,0.69}{\textbf{\textit{#1}}}}
\newcommand{\KeywordTok}[1]{\textcolor[rgb]{0.00,0.44,0.13}{\textbf{#1}}}
\newcommand{\NormalTok}[1]{#1}
\newcommand{\OperatorTok}[1]{\textcolor[rgb]{0.40,0.40,0.40}{#1}}
\newcommand{\OtherTok}[1]{\textcolor[rgb]{0.00,0.44,0.13}{#1}}
\newcommand{\PreprocessorTok}[1]{\textcolor[rgb]{0.74,0.48,0.00}{#1}}
\newcommand{\RegionMarkerTok}[1]{#1}
\newcommand{\SpecialCharTok}[1]{\textcolor[rgb]{0.25,0.44,0.63}{#1}}
\newcommand{\SpecialStringTok}[1]{\textcolor[rgb]{0.73,0.40,0.53}{#1}}
\newcommand{\StringTok}[1]{\textcolor[rgb]{0.25,0.44,0.63}{#1}}
\newcommand{\VariableTok}[1]{\textcolor[rgb]{0.10,0.09,0.49}{#1}}
\newcommand{\VerbatimStringTok}[1]{\textcolor[rgb]{0.25,0.44,0.63}{#1}}
\newcommand{\WarningTok}[1]{\textcolor[rgb]{0.38,0.63,0.69}{\textbf{\textit{#1}}}}

\usepackage{newunicodechar}
\newunicodechar{µ}{\ifmmode\mu\else\textmu\fi}
\newunicodechar{°}{\ifmmode^\circ\else\textdegree\fi}
\newunicodechar{·}{\ifmmode\cdot\else\textperiodcentered\fi}
\newunicodechar{²}{\ensuremath{^{2}}}
\newunicodechar{³}{\ensuremath{^{3}}}
\newunicodechar{¹}{\ensuremath{^{1}}}
\newunicodechar{⁰}{\ensuremath{^{0}}}
\newunicodechar{⁴}{\ensuremath{^{4}}}
\newunicodechar{⁵}{\ensuremath{^{5}}}
\newunicodechar{⁶}{\ensuremath{^{6}}}
\newunicodechar{⁷}{\ensuremath{^{7}}}
\newunicodechar{⁸}{\ensuremath{^{8}}}
\newunicodechar{⁹}{\ensuremath{^{9}}}
\newunicodechar{⁺}{\ensuremath{^{+}}}
\newunicodechar{⁻}{\ensuremath{^{-}}}
\newunicodechar{₀}{\ensuremath{_{0}}}
\newunicodechar{₁}{\ensuremath{_{1}}}
\newunicodechar{₂}{\ensuremath{_{2}}}
\newunicodechar{₃}{\ensuremath{_{3}}}
\newunicodechar{₄}{\ensuremath{_{4}}}
\newunicodechar{ₙ}{\ensuremath{_{n}}}
\newunicodechar{ν}{\ensuremath{\nu}}
\newunicodechar{Δ}{\ensuremath{\Delta}}
\newunicodechar{Ω}{\ensuremath{\Omega}}
\newunicodechar{α}{\ensuremath{\alpha}}
\newunicodechar{β}{\ensuremath{\beta}}
\newunicodechar{δ}{\ensuremath{\delta}}
\newunicodechar{ε}{\ensuremath{\varepsilon}}
\newunicodechar{θ}{\ensuremath{\theta}}
\newunicodechar{κ}{\ensuremath{\kappa}}
\newunicodechar{λ}{\ensuremath{\lambda}}
\newunicodechar{π}{\ensuremath{\pi}}
\newunicodechar{ρ}{\ensuremath{\rho}}
\newunicodechar{σ}{\ensuremath{\sigma}}
\newunicodechar{τ}{\ensuremath{\tau}}
\newunicodechar{→}{\ensuremath{\rightarrow}}
\newunicodechar{∎}{\ensuremath{\blacksquare}}
\newunicodechar{−}{\ensuremath{-}}
\newunicodechar{∧}{\ensuremath{\wedge}}
\newunicodechar{∫}{\ensuremath{\int}}
\newunicodechar{≈}{\ensuremath{\approx}}
\newunicodechar{≤}{\ensuremath{\leq}}
\newunicodechar{≥}{\ensuremath{\geq}}
\newunicodechar{⊆}{\ensuremath{\subseteq}}
\newunicodechar{∞}{\ensuremath{\infty}}
\newunicodechar{š}{\ifmmode\check{s}\else\v{s}\fi}
\newunicodechar{✓}{\ensuremath{\checkmark}}
\newunicodechar{≠}{\ensuremath{\neq}}
\newunicodechar{↔}{\ensuremath{\leftrightarrow}}
\newunicodechar{≪}{\ensuremath{\ll}}
\newunicodechar{↛}{\ensuremath{\nrightarrow}}
\newunicodechar{⇒}{\ensuremath{\Rightarrow}}
\newunicodechar{∪}{\ensuremath{\cup}}
\newunicodechar{⟨}{\ensuremath{\langle}}
\newunicodechar{⟩}{\ensuremath{\rangle}}
\newunicodechar{・}{\textperiodcentered}

\usepackage{xurl}
\usepackage{bookmark}
\usepackage{hyperref}

\setcounter{secnumdepth}{-\maxdimen}
\providecommand{\tightlist}{\setlength{\itemsep}{1pt}\setlength{\parskip}{0pt}}

\definecolor{openieAmber}{HTML}{D97706}
\definecolor{openieAmberBg}{HTML}{FFFBEB}
\definecolor{openieAmberBorder}{HTML}{F59E0B}
\definecolor{openieSlate}{HTML}{334155}
\definecolor{openieDark}{HTML}{0F172A}
\definecolor{openieCodeBg}{HTML}{F8FAFC}
\definecolor{openieCodeBorder}{HTML}{CBD5E1}
\definecolor{openieBlue}{HTML}{0057B8}

\hypersetup{
  pdftitle={Writing a Game Once and Running It Everywhere: Where Exact Replay Ends and Tolerance Begins},
  pdfauthor={David Charlot, Open Interface Engineering},
  colorlinks=true,
  linkcolor=openieBlue,
  citecolor=openieBlue,
  urlcolor=openieBlue,
  pdfcreator={XeTeX / Tectonic — Open Interface Engineering}
}

\urlstyle{same}
\setlength{\emergencystretch}{3em}
\hfuzz=6pt
\sloppy

\newcommand{\openiefig}[3]{%
  \IfFileExists{#1}{%
    \includegraphics[width=#2\linewidth,height=0.70\textheight,keepaspectratio]{#1}%
  }{%
    \fcolorbox{openieCodeBorder}{openieCodeBg}{%
      \parbox{0.92\linewidth}{\footnotesize\ttfamily Figure asset embedded in archival PDF; interactive/vector source: \url{#3}}%
    }%
  }%
}

\title{\vspace{-1.2em}\textbf{Writing a Game Once and Running It Everywhere: Where Exact Replay Ends and Tolerance Begins}}
\author{\textbf{David Charlot, Open Interface Engineering}\\[0.15em]
{\small Research study $\cdot$ Open Interface Engineering $\cdot$ \href{https://research.openie.dev}{research.openie.dev}}}
\date{\small Revision \texttt{1561a2a} (22 November 2026) $\cdot$ Web: \href{https://research.openie.dev/papers/game-engines/}{research.openie.dev/papers/game-engines/} $\cdot$ Archival LaTeX: \href{https://research.openie.dev/pdfs/latex/game-engines.pdf}{/pdfs/latex/game-engines.pdf} (\href{https://research.openie.dev/pdfs/latex/game-engines.tex}{.tex})}

\begin{document}
\maketitle
\vspace{-0.6em}

% OpenIE Measurement Honesty Callout Box
\noindent\fcolorbox{openieAmberBorder}{openieAmberBg}{%
  \parbox{\dimexpr\linewidth-2\fboxsep-2\fboxrule\relax}{%
    \small\color{openieDark}%
    \textbf{Measurement.} Research study, draft, not a journal final. \textbf{\texttt{GE-H1}} (\texttt{30} target $\times$ host cells across \texttt{6} physical machines, \texttt{3} OSs [\texttt{macOS 27.2.1}, \texttt{Linux 6.8/7.0}, \texttt{Windows 11 Pro 26300}], \texttt{3} ISAs, and \texttt{7} language/compiler projections, \texttt{100,000} \texttt{Q32.32} fixed-point ticks, \textbf{\texttt{0 / 25,500}} \texttt{SHA-256} checkpoint mismatches), \textbf{\texttt{GE-H2}} (\texttt{25,000}-tick 3-OS and cross-ISA \texttt{f64} drift horizons: \texttt{Windows UCRT} vs.\ \texttt{glibc} tick~\texttt{10}, \texttt{Apple libSystem} vs.\ \texttt{glibc} tick~\texttt{10}, \texttt{Windows UCRT} vs.\ \texttt{Apple libSystem} tick~\texttt{37}, \texttt{MSVC /fp:precise} vs.\ \texttt{Clang} \texttt{FMA} tick~\texttt{192}, cross-ISA \texttt{FMA} tick~\texttt{406}), and \textbf{\texttt{GE-H3}} (\texttt{8} hardware \texttt{GPU $\times$ API} endpoints across \texttt{DirectX 12}, \texttt{Metal}, and \texttt{Vulkan}: \textbf{\texttt{28 / 28}} hardware pairs pass at \texttt{PSNR >= 64.16 dB}, \texttt{SSIM >= 0.999833}, max $|\Delta| \le 2$, isolating same-GPU \texttt{DX12} vs.\ \texttt{Vulkan} shader-compiler drift at \texttt{79.52--108.34 dB}) were measured on physical hosts using \texttt{L2}/\texttt{L3} \texttt{reported\_j} and Win32 cycle counters. \texttt{GE-H4} to \texttt{GE-H9} remain open. No hardware rail or wall meter was attached (\texttt{measured\_j = None}, \texttt{board\_metered = false}).
  }%
}

\vspace{6pt}
% OpenIE Amber Left-Rule Summary Quote
\noindent{\color{openieAmber}\vrule width 3pt}\hspace{8pt}%
\parbox{\dimexpr\linewidth-11pt\relax}{%
  \small\itshape\color{openieSlate}%
  Can a game be written once as a checked state machine and projected into several engines, instruction sets, operating systems, and graphics stacks? Executed across 30 multi-ISA target × host cells for 100,000 ticks under Q32.32 integer physics (GE-H1: 0 / 25,500 state-digest mismatches across arm64 macOS, aarch64 Linux, Zen 5 and Raptor Lake x86\_64 Linux, and Tiger Lake x86\_64 Windows 11), 25,000 cross-platform floating-point ticks (GE-H2: t\_div = 10 Windows Universal C Runtime and Apple libm vs Linux glibc sin/cos, t\_div = 192–406 fused multiply-add lowering without -ffast-math), and 8 hardware graphics endpoints across Metal, Vulkan, and DirectX 12 (GE-H3: 0.00\%–9.34\% pixel drift, Δ\_max ≤ 2, peak signal-to-noise ratio ≥ 64.2 dB, structural similarity ≥ 0.99983 across all 28 pairs), paired with a playable joule-budgeted Q32.32 prototype.
}
\vspace{6pt}

{
\small
\setcounter{tocdepth}{2}
\tableofcontents
}
\vspace{4pt}
\hrule
\vspace{8pt}

\emph{Revision
\href{https://github.com/openIE-dev/research-openie-web/commits/main/src/content/papers/game-engines.md}{\texttt{1561a2a}},
22 November 2026. Interactive explorer, 3-OS / 3-API cross-GPU
inspector, and playable \texttt{Q32.32} Chrono-Budget Arena:
\href{https://research.openie.dev/living/game-engines/}{research.openie.dev/living/game-engines/}.}

\section{Abstract}\label{abstract}

This study asks whether a game can be written once as a checked state
machine and then projected across languages, each operating system (OS),
each central processing unit (CPU) instruction set architecture (ISA),
and each graphics processing unit (GPU) application programming
interface (API). Both the surveyed specifications and three executed
empirical falsifier suites give a sharp boundary: \textbf{yes} for game
rules and simulation under pinned fixed-point integer arithmetic
(\texttt{(GE-H1)}: \texttt{0\ /\allowbreak{}\ 25,\allowbreak{}500} Secure Hash Algorithm 256-bit
(\texttt{SHA-256}) checkpoint mismatches across \texttt{100,000} ticks
of a 16-body 2D arena projected across \texttt{30} target × host cells
spanning three operating systems and three instruction sets ---
\texttt{arm64} \texttt{macOS\ 27.\allowbreak{}2.\allowbreak{}1}, \texttt{aarch64}
\texttt{Linux\ 6.\allowbreak{}8}, Advanced Micro Devices (\texttt{AMD}) Zen 5 and
Intel Raptor Lake \texttt{x86\_64} \texttt{Linux\ 7.\allowbreak{}0}, and bare-metal
Intel Tiger Lake \texttt{x86\_64} \texttt{Windows\ 11} build
\texttt{26300} including Microsoft Visual C++ (\texttt{MSVC})
\texttt{cl.exe} --- spanning \texttt{0.073\ J} to \texttt{13.60\ J} in
macOS per-process \texttt{reported\_\allowbreak{}j} and \texttt{1.89\ J} to
\texttt{181.87\ J} in Linux Running Average Power Limit (\texttt{RAPL})
package \texttt{reported\_\allowbreak{}j}); \textbf{no} when unpinned floating-point
transcendentals or fused multiply-add (\texttt{FMA}) lowering enter
chaotic multi-body collisions (\texttt{(GE-H2)}: even \emph{without}
\texttt{-\allowbreak{}ffast-\allowbreak{}math} or \texttt{/fp:fast}, all three desktop OS math
libraries form a pairwise-divergent triangle --- Windows Universal C
Runtime (\texttt{UCRT}) vs Linux system \texttt{glibc}
\texttt{std::sin}/\allowbreak{}\texttt{std::cos} diverges on the same
\texttt{x86\_64} ISA at tick \texttt{10}, Apple \texttt{libSystem}
\texttt{libm} vs Linux \texttt{glibc} diverges at tick \texttt{10}, and
Windows \texttt{UCRT} vs Apple \texttt{libm} diverges at tick
\texttt{37} --- while \texttt{MSVC} \texttt{/\allowbreak{}fp:precise} vs
\texttt{Clang} \texttt{-\allowbreak{}ffp-\allowbreak{}contract=fast} diverges on \texttt{x86\_64}
at tick \texttt{192} and \texttt{x86\_64} \texttt{-\allowbreak{}mavx2\ -\allowbreak{}mfma} vs
\texttt{arm64} \texttt{fmadd} diverges at tick \texttt{406}); and
\textbf{no} for bit-exact pixels across GPU back ends, yet \textbf{yes}
under a bounded one-way tolerance gate (\texttt{(GE-H3)}: across eight
hardware \texttt{GPU\ ×\ API} endpoints spanning Apple \texttt{Metal},
Linux Mesa Radeon Vulkan (\texttt{RADV}), Nvidia (\texttt{NVIDIA}) Linux
and Windows \texttt{Vulkan} and Windows DirectX 12 (\texttt{DX12}), and
Intel Linux Mesa Vulkan (\texttt{ANV}) and Windows \texttt{Vulkan} and
\texttt{DX12} at \texttt{512×512} \texttt{RGBA8Unorm},
\texttt{27\ /\ 28} hardware pairs --- including the \textbf{same
physical GPU} across \texttt{DX12} vs \texttt{Vulkan} on both
\texttt{GeForce\ 2060} and \texttt{Iris\ Xe} --- differ in raw
\texttt{SHA-256}, yet all \texttt{28\ /\ 28} hardware pairs keep maximum
per-channel error \(\Delta_{\max} \le 2 / 255\), peak signal-to-noise
ratio (\texttt{PSNR}) \texttt{\textgreater{}=\ 64.16\ dB}, and
structural similarity (\texttt{SSIM}) \texttt{\textgreater{}=\ 0.99983},
with \texttt{measured\_\allowbreak{}j\ =\ None} and
\texttt{board\_\allowbreak{}metered\ =\ false}).

\textbf{In plain words}

\begin{enumerate}
\def\labelenumi{\arabic{enumi}.}
\tightlist
\item
  \textbf{Exact simulation across operating systems, chips, and
  languages works when arithmetic is integer-pinned.} Over
  \texttt{100,000} ticks (\texttt{25,500} verified \texttt{SHA-256}
  checkpoints across \texttt{30} cells), native Rust, C++20
  (\texttt{Clang}, the compiler collection (\texttt{GCC}), and Windows
  \texttt{MSVC}), WebAssembly (\texttt{Wasm}), JavaScript
  \texttt{BigInt}, and Python replay the exact same 732-byte state
  bit-for-bit across \texttt{macOS} (\texttt{arm64}), \texttt{Linux}
  (\texttt{aarch64} and \texttt{x86\_64}), and \texttt{Windows\ 11}
  (\texttt{x86\_64}) with zero mismatches.
\item
  \textbf{Standard floating-point physics drifts across
  \texttt{Windows}, \texttt{macOS}, and \texttt{Linux} even without
  fast-math flags.} Pure algebraic \texttt{+}, \texttt{-}, \texttt{*},
  \texttt{/}, and \texttt{sqrt} with \texttt{FMA} disabled
  (\texttt{/\allowbreak{}fp:strict} and \texttt{-\allowbreak{}ffp-\allowbreak{}contract=off}) match across
  every tested OS and ISA for all \texttt{25,000} ticks, but calling
  \texttt{std::sin}/\allowbreak{}\texttt{std::cos} across Windows \texttt{UCRT},
  macOS \texttt{libm}, and Linux \texttt{glibc} flips a bit within
  \texttt{10} to \texttt{37} ticks (\texttt{0.\allowbreak{}17–0.\allowbreak{}62\ s} at
  \texttt{60\ Hz}) and desynchronizes gameplay scores by tick
  \texttt{9,798}--\texttt{10,392} (\texttt{163.\allowbreak{}3–173.\allowbreak{}2\ s}), while
  default compiler \texttt{FMA} rules diverge between \texttt{MSVC}
  (\texttt{/\allowbreak{}fp:precise}) and \texttt{Clang}
  (\texttt{-\allowbreak{}ffp-\allowbreak{}contract=fast}) at tick \texttt{192} and between
  \texttt{x86\_64} and \texttt{arm64} at tick \texttt{406}.
\item
  \textbf{Even the same graphics card can flip pixel bits between
  \texttt{DirectX\ 12} and \texttt{Vulkan}, so rendering belongs
  downstream of a one-way tolerance gate.} Across eight hardware
  \texttt{GPU\ ×\ API} endpoints (\texttt{Metal}, \texttt{Vulkan}, and
  \texttt{DirectX\ 12}) plus Windows Advanced Rasterization Platform
  (\texttt{WARP}) and \texttt{Mesa\ llvmpipe}, switching
  \texttt{DirectX\ 12} vs \texttt{Vulkan} on the exact same
  \texttt{NVIDIA\ GeForce\ 2060} (\texttt{108.\allowbreak{}34\ dB} \texttt{PSNR}) or
  \texttt{Intel\ Iris\ Xe} (\texttt{79.52\ dB} \texttt{PSNR}) alters
  \texttt{1} least significant bit (\texttt{LSB}) via shader-compiler
  lowering, while cross-vendor pairs differ on \texttt{0.00\%} to
  \texttt{9.34\%} of pixels by \texttt{\textless{}=\ 2} \texttt{LSB}
  (\texttt{PSNR\ \textgreater{}=\ 64.16\ dB},
  \texttt{SSIM\ \textgreater{}=\ 0.99983}) without ever corrupting
  simulation state.
\end{enumerate}


\begin{center}
\openiefig{fig/game-engines-fig01-architecture.png}{0.94}{https://research.openie.dev/figures/game-engines/architecture.svg}
\end{center}


\emph{Figure 1. Proposed split. One Form feeds four projections.
Simulation and scripting sit in the exact zone, judged by per-tick state
hashes. Scene and rendering sit in the tolerance zone, judged by
per-pixel error. Host effects enter only as logged events.}


\begin{center}
\openiefig{fig/game-engines-fig02-coverage.png}{0.94}{https://research.openie.dev/figures/game-engines/coverage.svg}
\end{center}


\emph{Figure 2. Coverage of existing representations, scored only from
the documents read for this study. Green: documented capability. Amber:
partial. Grey: absent or undocumented. The rightmost column, cross-layer
checking, is nearly empty.}

\section{1. Terms, notation, and the idea under
test}\label{terms-notation-and-the-idea-under-test}

\begin{itemize}
\tightlist
\item
  RCE, IR, ECS, FSM, SDF: Reality Code Engine {[}same house{]},
  intermediate representation, entity component system, finite-state
  machine, and signed distance field
\item
  CPU, GPU, ISA, ISAs, SIMD, AVX2, NEON, SFU, ROP, RPL, TGL: central
  processing unit, graphics processing unit, instruction set
  architecture(s), single instruction multiple data, Advanced Vector
  Extensions 2, Arm Advanced SIMD, special function unit, raster
  operations pipeline, Intel Raptor Lake mobile architecture, and Intel
  Tiger Lake mobile architecture
\item
  IEEE, FMA, VFMADD, VFNMADD, FMADD, FMSUB, ULP, LSB, SHA, MAE, RMSE,
  PSNR, SSIM, SD, RGBA8: Institute of Electrical and Electronics
  Engineers (as in IEEE 754 floating-point), fused multiply-add, x86\_64
  vector fused multiply-add and negated multiply-add opcodes, AArch64
  scalar fused multiply-add and multiply-subtract opcodes, unit in the
  last place, least significant bit, Secure Hash Algorithm
  (\texttt{SHA-256}), mean absolute error, root mean squared error, peak
  signal-to-noise ratio, structural similarity index, standard
  deviation, and 8-bit-per-channel red-green-blue-alpha framebuffer
  format
\item
  WGSL, WESL, SPIR-V, GLSL, HLSL, MSL, D3D11, D3D12, DX12, VULKAN,
  METAL, DXIL, DXBC, DXC, FXC, WARP, WDDM, CUDA: WebGPU Shading
  Language, WebGPU Extended Shading Language, Standard Portable
  Intermediate Representation V, OpenGL Shading Language, High-Level
  Shading Language, Metal Shading Language, Direct3D 11, Direct3D 12,
  DirectX 12, \texttt{wgpu::Backends::VULKAN} and
  \texttt{wgpu::Backends::METAL} back-end flags, DirectX Intermediate
  Language, DirectX Bytecode, DirectX Shader Compiler, legacy Direct3D
  Effect-Compiler Tool, Windows Advanced Rasterization Platform software
  DirectX 12 rasterizer, Windows Display Driver Model, and Compute
  Unified Device Architecture
\item
  RADV, ANV, LLVM, GNU, GCC, MSVC, UCRT, NT, PE, OS, OSs, API, APIs,
  JIT, WebXR, WebMCP, URL, JSON, YAML, README, MB, CS0246, RE, ENGINE,
  DFU: Mesa Radeon Vulkan driver, Mesa Intel Vulkan driver, Low Level
  Virtual Machine compiler infrastructure, GNU's Not Unix project, GNU
  Compiler Collection, Microsoft Visual C++ compiler (\texttt{cl.exe}),
  Windows Universal C Runtime (\texttt{ucrtbase.\allowbreak{}dll}), Windows New
  Technology kernel (\texttt{Windows\ NT}), Windows Portable Executable,
  operating system(s), application programming interface(s),
  just-in-time compiler, Web Extended Reality API, Web Model Context
  Protocol, Uniform Resource Locator, JavaScript Object Notation, YAML
  Ain't Markup Language, repository README documentation file,
  megabytes, C\# compiler missing-type diagnostic CS0246, Capcom RE
  ENGINE (``Reach for the Moon'' / ``RE neXt'' engine branding), and
  Dagstuhl Follow-Ups DOI prefix
\item
  SCXML, GDL, VGDL, BSN, BSL, MLIR, ACID, GGPO, glTF, OpenUSD, USD, LAN,
  AI, CIG: State Chart XML, Game Description Language, Video Game
  Description Language, Bevy Scene Notation, Boldsea Semantic Language,
  Multi-Level Intermediate Representation,
  atomic/consistent/isolated/durable transactions, Good Game Peace Out
  rollback networking library, Graphics Language Transmission Format
  (ISO/IEC 12113:2022), Open Universal Scene Description, Universal
  Scene Description, local area network, artificial intelligence, and
  Computational Intelligence and Games
\item
  SATB, RAPL, MSR, PID, RUSAGE, V6, AMD, NVIDIA, RTX, V100, GH200, CTC:
  Satiation Tiered Benchmark telemetry levels (\texttt{L2}
  socket/package counter, \texttt{L3} per-process kernel counter)
  {[}same house{]}, Running Average Power Limit, model-specific
  register, process identifier, \texttt{proc\_\allowbreak{}pid\_\allowbreak{}rusage}
  (\texttt{RUSAGE\_\allowbreak{}INFO\_\allowbreak{}V6}), Advanced Micro Devices, NVIDIA
  Corporation, NVIDIA GeForce RTX GPU family, Volta V100 GPU, Grace
  Hopper GH200 superchip, and closed timelike curve
\end{itemize}

To keep forward references transparent from top to bottom, Section 1
defines the Form notation, Section 2 surveys existing game
representations, Section 3 analyzes floating-point divergence, Section 4
tabulates intermediate representations, Section 5 states the proposed
layer split, Section 6 reports the 3-OS multi-ISA and 3-API cross-GPU
falsifier runs, Section 7 lists infeasibility flags, and Section 8
documents sources and limits.

The \href{https://research.openie.dev/about/rce/}{Reality Code Engine
(RCE)} treats a program as a state machine called a Form, written
\(P = (S, E, \delta, D, \lambda)\):

\begin{itemize}
\tightlist
\item
  \(S\): the set of states;
\item
  \(E\): the set of events (inputs);
\item
  \(\delta : S \times E \to S\): the transition function;
\item
  \(D\): an append-only log of every step taken;
\item
  \(\lambda\): the output map from state to observable output.
\end{itemize}

A strict checker either accepts a Form or names the rule it breaks.
Emitters then project the accepted Form into Rust, Python, JavaScript
and Scala. XState (a JavaScript and TypeScript statechart library),
SCXML (a World Wide Web Consortium Recommendation) and WebAssembly
(\texttt{Wasm}) are not among the emitters' targets; this study treats
them as candidate projections. Because every step lands in \(D\), any
run can be replayed.

A game engine is a much larger machine. It runs a fixed-step simulation,
a scripting layer, a scene description, and a rendering pipeline on a
GPU. Sections 1 to 8 ask which of those layers can live inside a Form,
which must stay outside as host effects, and where exactness has to give
way to tolerance.

The idea has a long pedigree. The Stanford
\href{http://ggp.stanford.edu/readings/gdl_spec.pdf}{Game Description
Language (GDL) specification} covers ``finite, discrete, deterministic
multi-player games of complete information'' and observes that the
natural underlying model of such games ``is a finite state machine.''
What is new in a game engine is not the state machine but everything
around it: floating-point physics, scripts, scenes, and GPUs.

\section{2. What the landscape shows}\label{what-the-landscape-shows}

\subsection{2.1 Game logic as data already works, on one
backend}\label{game-logic-as-data-already-works-on-one-backend}

GameRule (\href{https://doi.org/10.1007/s00530-026-02629-8}{Multimedia
Systems 32:564, 2026}) expresses behavior as per-actor \emph{when-do}
rules over 14 create, read, update and delete actions and 6 conditions.
Its deterministic contract is concrete: one central scheduler runs from
a single FixedUpdate in declaration order, the last write wins, and
spawns and deletes are deferred to the end of the traversal. The authors
justify the contract by quoting Unity's manual, which says script
execution order ``isn't guaranteed across builds, machines, or Unity
versions.'' Across three Unity sample games (Haunted Jaunt, Survival
Shooter, Tanks), behavior code shrank from 572 to 107 lines (−81.3\%).

The price is measured. Rule interpretation costs about
\(c \approx 65\ \mu\text{s}\) per actor per step. With a step budget
\(T = 20\) ms, the actor ceiling is

\[N_{\max} = \frac{T}{c} = \frac{20\ \text{ms}}{65\ \mu\text{s}} \approx 308,\]

which agrees with the paper's observation that the budget is exceeded
between 200 and 350 actors (23.872 ms at 350). Only a Unity backend
exists, so portability across engines is argued but not demonstrated.

\subsection{2.2 Order of execution can be made deterministic by
construction}\label{order-of-execution-can-be-made-deterministic-by-construction}

The ECS pattern stores game state as components attached to entity
identifiers and runs behavior as systems over them. Redmond, Castello,
Calderón Trilla and Kuper give it a formal model, Core ECS
(\href{https://arxiv.org/abs/2508.15264}{arXiv:2508.15264}). They prove
that a schedule in which every concurrent pair of mutations commutes is
deterministic, and that two sub-schedules writing disjoint sets of
component labels may run in parallel without changing the result.
Surveying Bevy ECS, EnTT, Flecs, Specs and apecs, they find that none
allows attaching a new component concurrently with another mutation;
such structural changes are deferred and applied serially. Tasnim and
Zhao
(\href{https://doi.org/10.1145/3748522.3779910\%5D(https://doi.org/10.1145/3748522.3779910)}{SAC
2026,
{[}doi:10.1145/3748522.3779910{]}(https://doi.org/10.1145/3748522.3779910))})
add an operational semantics and a type system for archetype ECS that
tracks read and write sets per system, and report on a Tower Defense
simulation that a parallel archetype layout reached 109.54 frames per
second against 36.13 for an object-oriented baseline. GameRule's
deferred spawn and delete is the same discipline, reached independently.

\subsection{2.3 Compiling is not the same as being
right}\label{compiling-is-not-the-same-as-being-right}

The Mage benchmark
(\href{https://arxiv.org/abs/2605.07342}{arXiv:2605.07342}) ran 858
generation attempts over 26 Unity mechanic patterns. Without an IR, 43\%
of outputs passed at runtime, yet mechanism F1 (agreement between
requested and implemented mechanics) was only about 0.12. Conditioning
on an IR dropped runtime pass to 14\% to 21\% but lifted mechanism F1 to
0.82 to 1.00. Compile rate was anti-correlated with functional
correctness: the missing-type error CS0246 appeared 469 times under IR
conditioning and zero times without it. A behavior-only IR and a
full-scene IR were statistically indistinguishable (McNemar p = 1.0).
The lesson for RCE: a checker that only proves ``it compiles'' measures
the wrong quantity. Interface checks must be behavioral.

\subsection{2.4 Public benchmarks with replayable traces
exist}\label{public-benchmarks-with-replayable-traces-exist}

GameCraft-Bench
(\href{https://arxiv.org/abs/2606.17861}{arXiv:2606.17861};
\href{https://github.com/FreedomIntelligence/gamecraft-bench}{repository})
offers 140 Godot 4 two-dimensional tasks in 15 families. Submissions
carry replayable input traces. The strongest agent in the paper scored
41.46\%, and the repository later reports 65.70\%. It does not cover
Unity, Unreal, 3D, multiplayer, or large physics scenes.
\href{https://www.summerengine.com/templates}{Summer Engine} publishes
17 open-source game templates (rally, platformer, tower defense,
metroidvania, survivors, first-person shooter, LAN multiplayer, voxel,
and others). The General Video Game AI framework and its VGDL
(\href{https://arxiv.org/abs/1802.10363}{arXiv:1802.10363}) supply 120
single-player and 60 two-player games defined by sprite, interaction,
termination, and level-mapping sets.

\subsection{2.5 Convergence across industry and
research}\label{convergence-across-industry-and-research}

Several efforts are converging on the same shape: declarative scene and
pipeline files, a sandboxed scripting layer, a modern shader language,
and an event log.

\begin{itemize}
\tightlist
\item
  \href{https://github.com/poro/nAIVE}{nAIVE} is built in Rust on
  \texttt{wgpu} with an ECS (hecs), YAML scenes, materials and render
  pipelines, Lua 5.4 with per-entity sandboxing, Slang shaders
  cross-compiled to WGSL, Rapier physics, an event bus with YAML schema
  validation, and a headless test runner. Its
  \href{https://naive.dev/}{site} shows a 92-line YAML file declaring a
  five-pass deferred pipeline. Its README also describes GPU compute
  simulation driving 50,000+ entities, a choice that places that
  simulation on the tolerance side of the boundary drawn next.
\item
  Capcom describes REX (``RE neXt Engine'') as a staged evolution of RE
  ENGINE involving about 200 engineers
  (\href{https://www.capcom.co.jp/ir/english/data/oar/2025/re-engine.html}{Capcom
  IR}). Press coverage
  (\href{https://www.pcgamer.com/gaming-industry/game-development/capcom-announces-plans-to-transform-its-re-engine-into-an-ai-generation-game-engine/}{PC
  Gamer}) reports six programs, among them Re:Dox, which translates
  different types of data into a common set of rules, and RE:Flows,
  which converts game mechanic ideas into a standardised programming
  language; it also reports that more than 2,000 people use RE Engine.
\item
  Bevy's BSN proposal
  (\href{https://github.com/bevyengine/bevy/discussions/14437}{discussion
  14437}) defines scenes as composable Patch, EntityPatch, and
  DynamicScene values with a \texttt{bsn!} macro and an asset format,
  while leaving event handlers unresolved.
\item
  \href{https://github.com/clockworklabs/SpacetimeDB}{SpacetimeDB} runs
  game logic as reducers that execute as ACID transactions over tables
  and replays a commit log, which is the RCE log \(D\) under another
  name.
\item
  Executable Ontologies
  (\href{https://doi.org/10.48550/arXiv.2601.07964}{doi:10.48550/arXiv.2601.07964})
  model game behavior in BSL over an append-only, immutable event graph.
  The paper recommends a hybrid in which physics, rendering and
  animation keep traditional high-frequency approaches while quests and
  decisions run on the event graph, joined through aggregated events.
  That is the same layer split proposed here.
\end{itemize}

\section{3. The physics of divergence: why exactness is
binary}\label{the-physics-of-divergence-why-exactness-is-binary}

Floating-point addition is not associative. On parallel hardware the
order of a reduction varies run to run. A study of floating-point
non-associativity on parallel hardware
(\href{https://arxiv.org/abs/2408.05148}{arXiv:2408.05148}) defines the
scalar variability

\[V_s = 1 - \left| \frac{f_{nd}}{f_d} \right|,\]

where \(f_d\) is a deterministic reduction and \(f_{nd}\) a
nondeterministic one. Permuting a sum of \(10^6\) normally distributed
double-precision values changed the result by up to about
\(4.26 \times 10^{-13}\), and the maximum \(|V_s|\) grew as
\(\sqrt{n}\). A deterministic tree reduction cost about 0.2\% relative
to atomic accumulation on an NVIDIA V100 and up to 7.8\% on a GH200. In
1,000 graph-network training runs, every run produced distinct weights.

A game turns such tiny differences into visible ones. Rigid-body contact
and many-body dynamics are typically chaotic. Under a local Lyapunov
exponent \(\Lambda\), a perturbation grows as

\[\delta(t) = \delta_0 \, e^{\Lambda t},\]

so the time until it reaches a gameplay-visible size \(\Delta\) is

\[t^{*} = \frac{1}{\Lambda} \ln\frac{\Delta}{\delta_0}.\]

\emph{Illustrative model:} take \(\delta_0 \approx 6 \times 10^{-8}\)
(one ULP of single precision, relative), \(\Delta = 1\) (an object on
the other side of a wall), and an assumed
\(\Lambda = 2\ \text{s}^{-1}\). Then
\(t^{*} \approx 0.5 \ln(1.7 \times 10^{7}) \approx 8.3\) s. The
logarithm is the key: shrinking \(\delta_0\) a millionfold buys only
about 6.9 more seconds. Approximate agreement therefore cannot rescue
simulation. Either the state is bitwise identical or it eventually
diverges --- a prediction tested directly across 25,000 multi-ISA and
3-OS collision ticks below (\texttt{GE-H2}).

Practice agrees.
\href{https://gafferongames.com/post/floating_point_determinism/}{Gaffer
On Games} records that cross-machine determinism requires the same
compiler and instruction set, that Battlezone 2 hit AMD versus Intel
differences in transcendental functions, and that MotoGP replays could
not be shared between Xbox and personal computers. It cites Markstein's
advice to forbid FMA contraction and reassociation. The
\href{https://rapier.rs/docs/user_guides/rust/determinism/}{Rapier
guide} shows a working path: cross-platform determinism holds only with
its \texttt{enhanced-\allowbreak{}determinism} feature (incompatible with
\texttt{simd8}), on IEEE 754-2008 targets (Wasm qualifies), using
software transcendentals.

For rendering, the \href{https://www.w3.org/TR/WGSL/}{WGSL
specification} states that ``WGSL sometimes permits several possible
behaviors for a given feature. This is a portability hazard.'' Its
floating-point section goes further: ``WGSL does not specify a rounding
mode''; ``An implementation may reassociate operations''; it ``may fuse
operations if the transformed expression is at least as accurate as the
original''; subnormal inputs and outputs of many operations ``may be
flushed to zero''; and the \texttt{fma} built-in ``may expand to an
ordinary multiply \ldots{} and an add.'' Pixels are thus specified to
differ across implementations. The
\href{https://registry.khronos.org/SPIR-V/specs/unified1/SPIRV.html}{SPIR-V
specification} offers a narrower tool: its NoContraction decoration
marks an arithmetic instruction that ``cannot be combined with another
instruction to form a single operation,'' and such operations ``are not
allowed to reassociate.'' That is a per-instruction guard on one back
end, not a portable guarantee for a whole shader.

\section{4. The IR table}\label{the-ir-table}

{\footnotesize\setlength{\tabcolsep}{3.8pt}\renewcommand{\arraystretch}{1.15}
\begin{longtable}[]{@{}
  >{\raggedright\arraybackslash}p{(\linewidth - 6\tabcolsep) * \real{0.1442}}
  >{\raggedright\arraybackslash}p{(\linewidth - 6\tabcolsep) * \real{0.1442}}
  >{\raggedright\arraybackslash}p{(\linewidth - 6\tabcolsep) * \real{0.3750}}
  >{\raggedright\arraybackslash}p{(\linewidth - 6\tabcolsep) * \real{0.3365}}@{}}
\toprule\noalign{}
\begin{minipage}[b]{\linewidth}\raggedright
\textbf{Representation}
\end{minipage} & \begin{minipage}[b]{\linewidth}\raggedright
\textbf{Layer}
\end{minipage} & \begin{minipage}[b]{\linewidth}\raggedright
\textbf{Strongest property (documented)}
\end{minipage} & \begin{minipage}[b]{\linewidth}\raggedright
\textbf{What it lacks for a Form}
\end{minipage} \\
\midrule\noalign{}
\endhead
\bottomrule\noalign{}
\endlastfoot
\href{https://shader-slang.org/}{Slang} & Shaders & Module IR; back ends
for D3D12, Vulkan, Metal, D3D11, OpenGL, CUDA, WebGPU, CPU; capability
system & No game state, no events, no exact-semantics guarantee for
floats \\
\href{https://github.com/gfx-rs/wgpu/tree/trunk/naga}{Naga} & Shaders &
Validated IR; WGSL, SPIR-V and GLSL in; SPIR-V, WGSL, MSL, HLSL and GLSL
out & Same; no link to a scene schema \\
\href{https://www.w3.org/TR/WGSL/}{WGSL} /
\href{https://wesl-lang.dev/}{WESL} & Shaders & Uniformity analysis;
WESL adds imports, conditional compilation, packages, \texttt{wgsl-test}
& Permits reassociation, fusion and flush-to-zero; no host-side types \\
\href{https://registry.khronos.org/SPIR-V/specs/unified1/SPIRV.html}{SPIR-V}
& Shader binary IR & NoContraction decoration forbids fusion and
reassociation per instruction & Per-instruction only; no game
semantics \\
\href{https://mlir.llvm.org/}{MLIR} & Compiler infrastructure &
Multi-level dialects for lowering & Not a game model; semantics left to
each dialect \\
\href{https://www.khronos.org/gltf/}{glTF 2.0} (ISO/IEC 12113:2022) &
Assets and scene & Interchange standard across tools & No behavior, no
execution order \\
\href{https://openusd.org/release/intro.html}{OpenUSD} & Scene &
Composition arcs, namespace paths & Self-described ``Not a rigging
system''; no simulation semantics \\
\href{https://github.com/bevyengine/bevy/discussions/14437}{BSN} & Scene
& Typed patches over components, versioned imports & Event handlers
unresolved; Bevy only \\
\href{https://www.w3.org/TR/scxml/}{SCXML} /
\href{https://stately.ai/docs/xstate}{XState} & Behavior & Microstep and
macrostep semantics, document-order priority, serial parallel regions &
Data model is platform-specific; no numeric discipline, no scene or
shader \\
\href{http://ggp.stanford.edu/readings/gdl_spec.pdf}{GDL} & Rules &
Logic description of finite, discrete, deterministic games & No
continuous physics, no rendering \\
\href{https://arxiv.org/abs/1802.10363}{VGDL} & Rules & Compact 2D grid
games, large public corpus & 2D grid only; no rendering or physics
fidelity \\
\href{https://arxiv.org/abs/2508.15264}{Core ECS} & Execution order &
Theorem giving a condition for deterministic concurrent schedules & A
formal model, not an engine format; no scene or shader \\
\href{https://doi.org/10.1007/s00530-026-02629-8}{GameRule} & Rules &
Deterministic scheduler contract, measured code savings & Unity only;
about 65 µs per actor; no formal checker \\
\href{https://github.com/poro/nAIVE}{nAIVE} & Whole engine & YAML scenes
and pipelines, Lua, Slang, event schema validation & No published formal
semantics or cross-layer type check \\
\href{https://github.com/clockworklabs/SpacetimeDB}{SpacetimeDB} &
Server logic & Reducers as transactions, commit-log replay & Server side
only; no client rendering or scene \\
\end{longtable}
}

\section{5. Proposed architecture}\label{proposed-architecture}

\subsection{5.1 Layer split}\label{layer-split}

\begin{enumerate}
\def\labelenumi{\arabic{enumi}.}
\tightlist
\item
  \textbf{Simulation Form (exact).} \(S\) holds fixed-point or strict
  IEEE 754 values, with FMA and reassociation disabled and
  transcendentals implemented in integer or pinned software routines.
  \(E\) holds tick events and logged inputs. \(\delta\) uses SCXML
  macrostep semantics: per tick, process events in queue order, resolve
  conflicts by document order, run parallel regions serially.
  Collections iterate in a declared order, following GameRule's contract
  (declaration order, last write wins, structural changes deferred to
  the end of the tick). Where parallelism is wanted, the checker admits
  only schedules that satisfy the Core ECS disjointness condition.
\item
  \textbf{Scripting projection (exact).} Designer scripts compile to
  guards and actions over \(S\), emitted as
  \href{https://rhai.rs/book/about/features.html}{Rhai} (sandboxed, with
  a ``Don't Panic'' guarantee, no-std and Wasm support), Lua, or
  \href{https://luau.org/}{Luau} (sandboxed by default, with gradual
  types). Scripts may only call \(\delta\) through declared actions.
  They cannot read clocks, random sources, or GPU results directly.
\item
  \textbf{Scene projection (tolerance downstream).} \(\lambda\) maps
  \(S\) to a BSN-like patch tree whose component schema is part of the
  Form. glTF and OpenUSD serve as export targets for assets.
\item
  \textbf{Rendering projection (tolerance).} Materials and passes are
  authored in Slang or WGSL, and emitted through Slang or Naga to
  \texttt{DirectX\ 12}, \texttt{Metal}, and \texttt{Vulkan} back ends.
  Rendering is a pure function of \(\lambda(s)\). No GPU result ever
  flows back into \(S\).
\end{enumerate}

\subsection{5.2 Inside the Form versus host
effects}\label{inside-the-form-versus-host-effects}

Inside: state, transitions, rule order, fixed-point physics for
gameplay-relevant bodies, the random generator with its seed in \(S\),
interface schemas. Outside, as logged events: device input, wall-clock
time, network packets, audio output, asset input and output, GPU draws,
and any third-party physics used for cosmetic effects. Replay feeds
\(D\) back through \(\delta\) rather than touching devices. Storage cost
is modest. At a tick rate \(f_{\text{tick}} = 60\) Hz and \(b = 16\)
bytes of input per tick, the log grows at
\(f_{\text{tick}} \cdot b = 960\) B/s, about 3.5 MB per hour.

\subsection{5.3 Determinism boundary}\label{determinism-boundary}

\begin{itemize}
\tightlist
\item
  \textbf{Exact zone criterion:} for targets A and B and every tick
  \(t\), \(H(s_t^A) = H(s_t^B)\), where \(H\) is a cryptographic hash
  over a canonical serialization of \(S\).
\item
  \textbf{Tolerance zone criterion:} for rendered frames \(I^A\) and
  \(I^B\) of the same \(s_t\), the fraction of pixels with per-channel
  error above \(\varepsilon\) must not exceed \(p_{\max}\), paired with
  bounded structural similarity (\texttt{SSIM}) and peak signal-to-noise
  ratio (\texttt{PSNR}) floors. Section 6.3 executes this calibration
  across eight hardware \texttt{GPU\ ×\ API} endpoints spanning
  \texttt{DirectX\ 12}, \texttt{Metal}, and \texttt{Vulkan}
  (\texttt{Apple\ Metal}, \texttt{AMD\ RADV\ Vulkan}, \texttt{NVIDIA}
  Linux/Windows \texttt{Vulkan} and Windows \texttt{DX12}, and
  \texttt{Intel} Linux \texttt{ANV\ Vulkan} and Windows \texttt{Vulkan}
  / \texttt{DX12}).
\item
  Rollback networking in the style of
  \href{https://github.com/pond3r/ggpo}{GGPO}, which relies on input
  prediction and speculative execution, is valid only on the exact side
  of this boundary.
\end{itemize}

\subsection{5.4 Cross-layer interface
checking}\label{cross-layer-interface-checking}

The checker validates four typed edges before any emission: (a) script
to simulation: every action names a declared transition with matching
argument types; (b) simulation to scene: the output of \(\lambda\)
matches the component schema; (c) scene to rendering: every material
binding matches a Slang or WGSL resource of compatible type, group, and
binding index; (d) event coverage: every event a script listens for is
producible by some transition or host source. Mage's finding (Section
2.3) is the reason edge checks are judged by behavior traces, not by
compile success.

\section{6. Falsifier-first test plan and executed 3-OS / 3-ISA / 3-API
results}\label{falsifier-first-test-plan-and-executed-3-os-3-isa-3-api-results}

Each hypothesis is listed with its kill criterion. If the criterion is
met, the claim is withdrawn. Subsections 6.1 to 6.4 report the three
executed empirical falsifier suites (\texttt{GE-H1} 3-OS × 3-ISA
\texttt{Q32.32} replay across \texttt{30} target × host cells,
\texttt{GE-H2} 3-OS × 3-ISA floating-point drift across
\texttt{Windows\ 11\ UCRT}, \texttt{macOS\ libSystem}, and
\texttt{Linux\ glibc}, and \texttt{GE-H3} 3-OS × 3-API headless
\texttt{wgpu} rendering tolerance across \texttt{DirectX\ 12},
\texttt{Metal}, and \texttt{Vulkan}), plus an interactive joule-budgeted
\texttt{Q32.32} prototype (\texttt{GE-H4} through \texttt{GE-H9} remain
open).

{\scriptsize\setlength{\tabcolsep}{3.0pt}\renewcommand{\arraystretch}{1.14}
\begin{longtable}[]{@{}
  >{\raggedright\arraybackslash}p{(\linewidth - 10\tabcolsep) * \real{0.0356}}
  >{\raggedright\arraybackslash}p{(\linewidth - 10\tabcolsep) * \real{0.1733}}
  >{\raggedright\arraybackslash}p{(\linewidth - 10\tabcolsep) * \real{0.2311}}
  >{\raggedright\arraybackslash}p{(\linewidth - 10\tabcolsep) * \real{0.2133}}
  >{\raggedright\arraybackslash}p{(\linewidth - 10\tabcolsep) * \real{0.1822}}
  >{\raggedright\arraybackslash}p{(\linewidth - 10\tabcolsep) * \real{0.1644}}@{}}
\toprule\noalign{}
\begin{minipage}[b]{\linewidth}\raggedright
\textbf{ID}
\end{minipage} & \begin{minipage}[b]{\linewidth}\raggedright
\textbf{Hypothesis}
\end{minipage} & \begin{minipage}[b]{\linewidth}\raggedright
\textbf{Test corpus}
\end{minipage} & \begin{minipage}[b]{\linewidth}\raggedright
\textbf{Measure}
\end{minipage} & \begin{minipage}[b]{\linewidth}\raggedright
\textbf{Kill criterion}
\end{minipage} & \begin{minipage}[b]{\linewidth}\raggedright
\textbf{Status}
\end{minipage} \\
\midrule\noalign{}
\endhead
\bottomrule\noalign{}
\endlastfoot
GE-H1 & Simulation projections are bitwise identical across languages,
OSs, and CPU ISAs & 16-body 2D arena Form (\texttt{Q32.32} fixed-point,
\texttt{732}-byte state) across \texttt{arm64} macOS, \texttt{aarch64}
Linux, AMD/Intel \texttt{x86\_64} Linux, and Tiger Lake \texttt{x86\_64}
Windows 11 & \(H(s_t)\) across Rust, Wasm32, C++20
(\texttt{Clang}/\allowbreak{}\texttt{GCC}/\allowbreak{}\texttt{MSVC}), Node.js \texttt{BigInt},
Python 3 for \(10^5\) ticks (\texttt{30} cells) & Any hash mismatch at
any checkpoint on any \texttt{(\allowbreak{}host,\allowbreak{}\ target)} cell & \textbf{Executed
(\texttt{0\ /\allowbreak{}\ 25,\allowbreak{}500} mismatches across \texttt{3} OSs \& \texttt{3}
ISAs, Section 6.1)} \\
GE-H2 & Strict mode vs 3-OS / cross-ISA floating-point boundary \&
library/FMA drift & 16-body 2D chaotic arena (\texttt{25,000} ticks
across \texttt{arm64} macOS, \texttt{aarch64} Linux, \texttt{x86\_64}
Linux, and \texttt{x86\_64} Windows 11) & Exact divergence ticks
\texttt{t\_div} across strict \texttt{f64}, 3-OS
\texttt{UCRT}/\allowbreak{}\texttt{libSystem}/\allowbreak{}\texttt{glibc} \texttt{sin/cos},
\texttt{MSVC} vs \texttt{Clang} \texttt{FMA}, V8 \texttt{Math.sin}, and
\texttt{f32} & Unexpected mismatch under pinned algebraic arithmetic, or
unpredicted silent drift & \textbf{Executed across \texttt{macOS},
\texttt{Linux}, and \texttt{Windows\ 11} (Section 6.2)} \\
GE-H3 & Cross-GPU \& cross-API rendering stays within bounded tolerance
(\texttt{wgpu} / WGSL) & 4 canonical \texttt{Q32.32} states
(\texttt{t\ =\ 0,\allowbreak{}\ 250,\allowbreak{}\ 1k,\allowbreak{}\ 10k}) rendered at \texttt{512×512}
\texttt{RGBA8Unorm} on \texttt{8} hardware \texttt{GPU\ ×\ API}
endpoints (\texttt{DX12}, \texttt{Metal}, \texttt{Vulkan}) +
\texttt{WARP} + \texttt{llvmpipe} + \texttt{CPU\ Q16.\allowbreak{}16} & Framebuffer
\texttt{SHA-256}, differing pixel \texttt{\%}, \(\Delta_{\max}\), MAE,
RMSE, \texttt{PSNR}, and \texttt{SSIM} & Any physical GPU pair failing
\texttt{PSNR\ \textgreater{}\ 40\ dB},
\texttt{SSIM\ \textgreater{}\ 0.999}, or \(p(\Delta > 8) \le 0.001\) &
\textbf{Executed (\texttt{28/28} pairs = \texttt{112/112} frames pass
gate: \texttt{PSNR\ ≥\ 64.\allowbreak{}16\ dB}, \texttt{SSIM\ ≥\ 0.\allowbreak{}99983}, Section
6.3)} \\
GE-H4 & Form-encoded rules cut code as GameRule did & Haunted Jaunt,
Survival Shooter, Tanks behaviors & Lines of behavior code & Reduction
under 50\% (GameRule reached 81.3\%) & Open \\
GE-H5 & Form runtime exceeds the GameRule ceiling & Synthetic actor
scaling & Step time against \(N\) & Step time over 20 ms below
\(N = 350\) & Open \\
GE-H6 & Interface checks catch real defects & Mage's 26 patterns;
injected binding and schema faults & Detection rate; mechanism F1 &
Detection below 90\% of injected faults, or F1 no better than the no-IR
0.12 & Open \\
GE-H7 & Full \(\ln(\Delta/\delta_0)\) slope sweep on Summer Engine
templates & Physics-heavy Summer Engine templates re-expressed as Forms
& Divergence time against injected \(\delta_0\) & Slope of \(t^{*}\)
against \(\ln(1/\delta_0)\) not distinguishable from zero & Open \\
GE-H8 & Replay is complete & One-hour recorded sessions & Final hash
reproduced from \(D\) alone & Any replay that needs state not in \(D\) &
Open \\
GE-H9 & Benchmark parity & GameCraft-Bench 140 tasks via a Form-to-Godot
emitter & Pass rate against native Godot submissions & Lower pass rate
on behavior-identical inputs & Open \\
\end{longtable}
}


\begin{center}
\openiefig{fig/game-engines-fig03-cross-gpu-tolerance.png}{0.94}{https://research.openie.dev/figures/game-engines/cross-gpu-tolerance.svg}
\end{center}


\emph{Figure 3. Empirical boundary across \texttt{GE-H1},
\texttt{GE-H2}, and \texttt{GE-H3}. Upstream \texttt{Q32.32} simulation
state (\texttt{S}) achieves \texttt{0\ /\allowbreak{}\ 25,\allowbreak{}500} \texttt{SHA-256}
checkpoint mismatches across six machines, three OSs (\texttt{macOS},
\texttt{Linux}, \texttt{Windows\ 11}), and three ISAs, whereas unpinned
\texttt{f64} \texttt{sin/cos} and \texttt{FMA} diverge across
OS/compiler/ISA boundaries within \texttt{10} to \texttt{406} ticks.
Downstream of \(\lambda(s_t)\), switching \texttt{DirectX\ 12} vs
\texttt{Vulkan} on the exact same Windows 11 GPU flips \texttt{1} LSB
via shader lowering (\texttt{79.\allowbreak{}5–108.\allowbreak{}3\ dB} \texttt{PSNR}), while all
\texttt{28\ /\ 28} hardware \texttt{GPU\ ×\ API} pairs stay strictly
above \texttt{64.16\ dB} \texttt{PSNR} and \texttt{0.99983}
\texttt{SSIM} (\texttt{measured\_\allowbreak{}j\ =\ None},
\texttt{board\_\allowbreak{}metered\ =\ false}).}

\subsection{\texorpdfstring{6.1 \texttt{GE-H1}: Three-OS and three-ISA
100,000-tick determinism (\texttt{Q32.32} integer
physics)}{6.1 GE-H1: Three-OS and three-ISA 100,000-tick determinism (Q32.32 integer physics)}}\label{ge-h1-three-os-and-three-isa-100000-tick-determinism-q32.32-integer-physics}

\emph{Pre-registered under \texttt{PREREG\_\allowbreak{}wave3\_\allowbreak{}multi\_\allowbreak{}isa\_\allowbreak{}wgpu.\allowbreak{}json}
(\texttt{SHA-\allowbreak{}256:\ 1721e5bdffaf.\allowbreak{}.\allowbreak{}.\allowbreak{}}). Telemetry strictly separates
\texttt{SATB} levels \texttt{L2} and \texttt{L3} \texttt{reported\_\allowbreak{}j}
(\texttt{L2}: Linux \texttt{/\allowbreak{}sys/\allowbreak{}class/\allowbreak{}powercap/\allowbreak{}intel-\allowbreak{}rapl:0/\allowbreak{}energy\_\allowbreak{}uj}
socket/package and core counters; \texttt{L3}: macOS per-PID
\texttt{proc\_\allowbreak{}pid\_\allowbreak{}rusage(\allowbreak{}child\_\allowbreak{}pid,\allowbreak{}\ RUSAGE\_\allowbreak{}INFO\_\allowbreak{}V6).\allowbreak{}ri\_\allowbreak{}energy\_\allowbreak{}nj},
plus Win32 \texttt{GetProcessTimes} and \texttt{QueryProcessCycleTime}
activity accounting on Windows 11 where unprivileged ring-3
\texttt{RAPL} \texttt{MSR} access is blocked;
\texttt{measured\_\allowbreak{}j\ =\ None}, \texttt{board\_\allowbreak{}metered\ =\ false}).
Interactive explorer and raw JSON:
\href{https://research.openie.dev/living/game-engines/}{\texttt{/\allowbreak{}living/\allowbreak{}game-\allowbreak{}engines/\allowbreak{}}}
(\href{https://research.openie.dev/living/game-engines/results.json}{\texttt{results.\allowbreak{}json}}).}

A 16-body + 4-pickup 2D elastic-collision and cooldown-FSM arena Form
(\texttt{Q32.32} signed 64-bit fixed-point coordinates in
\texttt{{[}-64,\ +64{]}}, \texttt{128}-bit intermediate products,
integer floor square root \texttt{isqrt\_\allowbreak{}u128}, deferred structural
pickup respawn, \texttt{732}-byte little-endian canonical state vector)
was executed for \texttt{100,000} ticks across \textbf{six machines,
three ISAs (\texttt{arm64}/\allowbreak{}\texttt{AArch64}, \texttt{x86\_64},
\texttt{wasm32}), three operating systems (\texttt{macOS\ 27.\allowbreak{}2.\allowbreak{}1},
\texttt{Ubuntu\ Linux\ 6.\allowbreak{}8\ /\allowbreak{}\ 7.\allowbreak{}0},
\texttt{Microsoft\ Windows\ 11\ Pro} build \texttt{26300}), and seven
language/compiler projections (\texttt{30} target × host cells)}:

{\scriptsize\setlength{\tabcolsep}{3.0pt}\renewcommand{\arraystretch}{1.14}
\begin{longtable}[]{@{}
  >{\raggedright\arraybackslash}p{(\linewidth - 12\tabcolsep) * \real{0.1871}}
  >{\raggedright\arraybackslash}p{(\linewidth - 12\tabcolsep) * \real{0.1754}}
  >{\raggedleft\arraybackslash}p{(\linewidth - 12\tabcolsep) * \real{0.1287}}
  >{\raggedleft\arraybackslash}p{(\linewidth - 12\tabcolsep) * \real{0.0994}}
  >{\raggedright\arraybackslash}p{(\linewidth - 12\tabcolsep) * \real{0.1637}}
  >{\raggedleft\arraybackslash}p{(\linewidth - 12\tabcolsep) * \real{0.1111}}
  >{\centering\arraybackslash}p{(\linewidth - 12\tabcolsep) * \real{0.1345}}@{}}
\toprule\noalign{}
\begin{minipage}[b]{\linewidth}\raggedright
\textbf{Host (\texttt{ISA} · \texttt{OS})}
\end{minipage} & \begin{minipage}[b]{\linewidth}\raggedright
\textbf{Target projection}
\end{minipage} & \begin{minipage}[b]{\linewidth}\raggedleft
\textbf{Wall ms / \texttt{100k} (mean ± SD)}
\end{minipage} & \begin{minipage}[b]{\linewidth}\raggedleft
\textbf{\texttt{ns\ /\allowbreak{}\ tick}}
\end{minipage} & \begin{minipage}[b]{\linewidth}\raggedright
\textbf{Telemetry channel (\texttt{reported\_\allowbreak{}j})}
\end{minipage} & \begin{minipage}[b]{\linewidth}\raggedleft
\textbf{\texttt{reported\_\allowbreak{}j} / \texttt{100k} (J)}
\end{minipage} & \begin{minipage}[b]{\linewidth}\centering
\textbf{Verified \texttt{SHA-256} checkpoints}
\end{minipage} \\
\midrule\noalign{}
\endhead
\bottomrule\noalign{}
\endlastfoot
\textbf{Apple M5 Max} (\texttt{arm64} · \texttt{macOS\ 27.\allowbreak{}2.\allowbreak{}1}) &
\texttt{rust\_\allowbreak{}native} (\texttt{rustc\ 1.\allowbreak{}100\ -\allowbreak{}O}) &
\texttt{22.\allowbreak{}333\ ±\ 5.\allowbreak{}043} & \texttt{223.3} & \texttt{L3} per-PID
\texttt{ri\_\allowbreak{}energy\_\allowbreak{}nj} & \textbf{\texttt{0.\allowbreak{}1028\ ±\ 0.\allowbreak{}0247}} &
\texttt{100\ /\allowbreak{}\ 100} (\texttt{stride\ 1,\allowbreak{}000}) \\
\textbf{Apple M5 Max} (\texttt{arm64} · \texttt{macOS\ 27.\allowbreak{}2.\allowbreak{}1}) &
\texttt{cpp20\_\allowbreak{}clang\_\allowbreak{}o3} (\texttt{Apple\ Clang\ 21}) &
\texttt{21.\allowbreak{}022\ ±\ 1.\allowbreak{}657} & \texttt{210.2} & \texttt{L3} per-PID
\texttt{ri\_\allowbreak{}energy\_\allowbreak{}nj} & \textbf{\texttt{0.\allowbreak{}1090\ ±\ 0.\allowbreak{}0134}} &
\texttt{100\ /\allowbreak{}\ 100} (\texttt{stride\ 1,\allowbreak{}000}) \\
\textbf{Apple M5 Max} (\texttt{arm64} · \texttt{macOS\ 27.\allowbreak{}2.\allowbreak{}1}) &
\texttt{rust\_\allowbreak{}wasm32\_\allowbreak{}node} (\texttt{Node\ v26.\allowbreak{}13}) &
\texttt{84.\allowbreak{}300\ ±\ 11.\allowbreak{}334} & \texttt{843.0} & \texttt{L3} per-PID
\texttt{ri\_\allowbreak{}energy\_\allowbreak{}nj} & \textbf{\texttt{0.\allowbreak{}4480\ ±\ 0.\allowbreak{}0881}} &
\texttt{100\ /\allowbreak{}\ 100} (\texttt{stride\ 1,\allowbreak{}000}) \\
\textbf{Apple M5 Max} (\texttt{arm64} · \texttt{macOS\ 27.\allowbreak{}2.\allowbreak{}1}) &
\texttt{javascript\_\allowbreak{}bigint\_\allowbreak{}node} (\texttt{ES2024}) &
\texttt{717.\allowbreak{}955\ ±\ 36.\allowbreak{}304} & \texttt{7,179.6} & \texttt{L3} per-PID
\texttt{ri\_\allowbreak{}energy\_\allowbreak{}nj} & \textbf{\texttt{4.\allowbreak{}4021\ ±\ 0.\allowbreak{}4048}} &
\texttt{100\ /\allowbreak{}\ 100} (\texttt{stride\ 1,\allowbreak{}000}) \\
\textbf{Apple M5 Max} (\texttt{arm64} · \texttt{macOS\ 27.\allowbreak{}2.\allowbreak{}1}) &
\texttt{python\_\allowbreak{}3\_\allowbreak{}13} (\texttt{CPython\ 3.\allowbreak{}13.\allowbreak{}14}) &
\texttt{2,\allowbreak{}945.\allowbreak{}816\ ±\ 284.\allowbreak{}033} & \texttt{29,458.2} & \texttt{L3} per-PID
\texttt{ri\_\allowbreak{}energy\_\allowbreak{}nj} & \textbf{\texttt{13.\allowbreak{}5999\ ±\ 1.\allowbreak{}9086}} &
\texttt{100\ /\allowbreak{}\ 100} (\texttt{stride\ 1,\allowbreak{}000}) \\
\textbf{Apple M3 Ultra} (\texttt{arm64} · \texttt{macOS\ 27.\allowbreak{}2.\allowbreak{}1}) &
\texttt{rust\_\allowbreak{}native} / \texttt{cpp20\_\allowbreak{}clang\_\allowbreak{}o3} & \texttt{32.711} /
\texttt{33.727} & \texttt{327.1} / \texttt{337.3} & \texttt{L3} per-PID
\texttt{ri\_\allowbreak{}energy\_\allowbreak{}nj} & \textbf{\texttt{0.0729} / \texttt{0.0755}} &
\texttt{2\ ×\ 1,\allowbreak{}000\ /\allowbreak{}\ 1,\allowbreak{}000} (\texttt{stride\ 100}) \\
\textbf{AMD Ryzen AI Max+ 395} (\texttt{Zen\ 5} \texttt{x86\_64} ·
\texttt{Linux\ 7.\allowbreak{}0}) & \texttt{rust\_\allowbreak{}native} / \texttt{cpp20\_\allowbreak{}clang\_\allowbreak{}o3}
/ \texttt{cpp20\_\allowbreak{}gcc\_\allowbreak{}o3} & \texttt{21.793} / \texttt{23.061} /
\texttt{26.677} & \texttt{217.9} -- \texttt{266.8} & \texttt{L2} AMD
RAPL pkg (core \texttt{0.13}--\texttt{0.21\ J}) & \textbf{\texttt{1.958}
/ \texttt{2.349} / \texttt{2.932}} & \texttt{3\ ×\ 1,\allowbreak{}000\ /\allowbreak{}\ 1,\allowbreak{}000}
(\texttt{stride\ 100}) \\
\textbf{AMD Ryzen AI Max+ 395} (\texttt{Zen\ 5} \texttt{x86\_64} ·
\texttt{Linux\ 7.\allowbreak{}0}) & \texttt{rust\_\allowbreak{}wasm32} / \texttt{js\_\allowbreak{}bigint} /
\texttt{python\_\allowbreak{}3\_\allowbreak{}14} & \texttt{145.29} / \texttt{1,347.80} /
\texttt{2,546.42} & \texttt{1,452.9} -- \texttt{25,464.2} & \texttt{L2}
AMD RAPL pkg (core \texttt{1.01}--\texttt{12.23\ J}) &
\textbf{\texttt{14.44} / \texttt{115.81} / \texttt{181.87}} &
\texttt{3\ ×\ 1,\allowbreak{}000\ /\allowbreak{}\ 1,\allowbreak{}000} (\texttt{stride\ 100}) \\
\textbf{Intel Core i9-13900H} (\texttt{Raptor\ Lake} \texttt{x86\_64} ·
\texttt{Linux\ 7.\allowbreak{}0}) & \texttt{rust\_\allowbreak{}native} / \texttt{cpp20\_\allowbreak{}clang\_\allowbreak{}o3}
/ \texttt{cpp20\_\allowbreak{}gcc\_\allowbreak{}o3} & \texttt{39.718} / \texttt{36.868} /
\texttt{58.792} & \texttt{368.7} -- \texttt{587.9} & \texttt{L2} Intel
RAPL pkg (core \texttt{1.64}--\texttt{2.68\ J}) & \textbf{\texttt{2.062}
/ \texttt{1.888} / \texttt{3.026}} & \texttt{3\ ×\ 1,\allowbreak{}000\ /\allowbreak{}\ 1,\allowbreak{}000}
(\texttt{stride\ 100}) \\
\textbf{Intel Core i9-13900H} (\texttt{Raptor\ Lake} \texttt{x86\_64} ·
\texttt{Linux\ 7.\allowbreak{}0}) & \texttt{rust\_\allowbreak{}wasm32} / \texttt{js\_\allowbreak{}bigint} /
\texttt{python\_\allowbreak{}3\_\allowbreak{}14} & \texttt{158.35} / \texttt{1,390.46} /
\texttt{3,375.07} & \texttt{1,583.5} -- \texttt{33,750.7} & \texttt{L2}
Intel RAPL pkg (core \texttt{5.65}--\texttt{124.90\ J}) &
\textbf{\texttt{6.89} / \texttt{60.79} / \texttt{149.48}} &
\texttt{3\ ×\ 1,\allowbreak{}000\ /\allowbreak{}\ 1,\allowbreak{}000} (\texttt{stride\ 100}) \\
\textbf{NVIDIA Jetson Thor} (\texttt{Armv9-A} \texttt{aarch64} ·
\texttt{Linux\ 6.\allowbreak{}8}) & \texttt{rust\_\allowbreak{}native} / \texttt{clang} /
\texttt{gcc} / \texttt{wasm32} / \texttt{js} / \texttt{py3.12} &
\texttt{39.28} (\texttt{GCC}) -- \texttt{4,407.89} (\texttt{Py}) &
\texttt{392.8} -- \texttt{44,078.9} & CPU-only
(\texttt{reported\_\allowbreak{}j\ =\ null}, GPU untouched) & \texttt{null} &
\texttt{6\ ×\ 1,\allowbreak{}000\ /\allowbreak{}\ 1,\allowbreak{}000} (\texttt{stride\ 100}) \\
\textbf{Intel Core i7-1165G7} (\texttt{Tiger\ Lake} \texttt{x86\_64} ·
\texttt{Windows\ 11\ 26300}) & \texttt{rust\_\allowbreak{}native}
(\texttt{x86\_\allowbreak{}64-\allowbreak{}pc-\allowbreak{}windows-\allowbreak{}msvc}) / \texttt{cpp20\_\allowbreak{}msvc\_\allowbreak{}o2}
(\texttt{cl.\allowbreak{}exe\ 19.\allowbreak{}44}) & \texttt{49.\allowbreak{}730\ ±\ 1.\allowbreak{}727} /
\texttt{343.\allowbreak{}363\ ±\ 3.\allowbreak{}827} & \texttt{497.3} / \texttt{3,433.6} & Win32
\texttt{L3} cycles (\texttt{138.9M} / \texttt{958.9M}; CPU \texttt{RAPL}
blocked) & \texttt{null} & \texttt{2\ ×\ 1,\allowbreak{}000\ /\allowbreak{}\ 1,\allowbreak{}000}
(\texttt{stride\ 100}) \\
\textbf{Intel Core i7-1165G7} (\texttt{Tiger\ Lake} \texttt{x86\_64} ·
\texttt{Windows\ 11\ 26300}) & \texttt{rust\_\allowbreak{}wasm32}
(\texttt{Node\ v26.\allowbreak{}13}) / \texttt{js\_\allowbreak{}bigint} / \texttt{python\_\allowbreak{}3\_\allowbreak{}10} &
\texttt{198.01} / \texttt{1,776.96} / \texttt{8,899.23} &
\texttt{1,980.1} -- \texttt{88,992.3} & Win32 \texttt{L3} cycles
(\texttt{553.8M} -- \texttt{24,722.2M}; CPU \texttt{RAPL} blocked) &
\texttt{null} & \texttt{3\ ×\ 1,\allowbreak{}000\ /\allowbreak{}\ 1,\allowbreak{}000} (\texttt{stride\ 100}) \\
\end{longtable}
}

\begin{itemize}
\tightlist
\item
  \textbf{Verdict}: \textbf{\texttt{PASS\_\allowbreak{}ZERO\_\allowbreak{}MISMATCHES}
  (\texttt{0\ /\allowbreak{}\ 25,\allowbreak{}500} verified \texttt{SHA-256} checkpoints across
  \texttt{30} cells; \texttt{0\ /\allowbreak{}\ 300,\allowbreak{}000} pairwise stride-\texttt{100}
  comparisons across the \texttt{25} full-trace cells spanning
  \texttt{macOS}, \texttt{Linux}, and \texttt{Windows\ 11})}. Every
  target on every host produced \texttt{2683c3f7cb8b.\allowbreak{}.\allowbreak{}.\allowbreak{}} at tick
  \texttt{100}, \texttt{21fa688e4ee8.\allowbreak{}.\allowbreak{}.\allowbreak{}} at tick \texttt{1,000},
  \texttt{ea01ba793841.\allowbreak{}.\allowbreak{}.\allowbreak{}} at tick \texttt{50,000}, and
  \texttt{82e404df\allowbreak{}6475d64d.\allowbreak{}.\allowbreak{}.\allowbreak{}} at tick \texttt{100,000}.
\item
  \textbf{Windows \texttt{MSVC}
  \texttt{\textless{}\_\_msvc\_int128.hpp\textgreater{}} vs
  \texttt{LLVM} \texttt{i128} codegen}: On bare-metal Windows 11
  (\texttt{x86\_64} \texttt{PE}), \texttt{MSVC}
  (\texttt{cl.\allowbreak{}exe\ 19.\allowbreak{}44.\allowbreak{}35207\ /\allowbreak{}std:c++20\ /\allowbreak{}O2\ /\allowbreak{}fp:strict}) lacks
  built-in \texttt{\_\allowbreak{}\_\allowbreak{}int128} scalar types and lowers
  \texttt{\textless{}\_\_msvc\_int128.hpp\textgreater{}}
  \texttt{std::\_\allowbreak{}Signed128} and \texttt{std::\_\allowbreak{}Unsigned128} arithmetic
  to multi-word library routines (\texttt{343.\allowbreak{}36\ ms}, \texttt{958.9M}
  CPU cycles per \texttt{100k} ticks --- \texttt{6.90×} slower than
  \texttt{rustc\ 1.\allowbreak{}100} LLVM native \texttt{x86\_\allowbreak{}64-\allowbreak{}pc-\allowbreak{}windows-\allowbreak{}msvc} at
  \texttt{49.73\ ms} / \texttt{138.9M} cycles, and \texttt{1.73×} slower
  than \texttt{Node.js} \texttt{Wasm32} JIT at \texttt{198.\allowbreak{}01\ ms}), yet
  preserves \texttt{100\%} bit-for-bit state parity across all
  \texttt{1,000} checkpoints.
\item
  \textbf{Telemetry scope across \texttt{macOS}, \texttt{Linux}, and
  \texttt{Windows\ 11}}: Notice why \texttt{L2} and \texttt{L3} channels
  must never be conflated: Apple \texttt{L3} \texttt{ri\_\allowbreak{}energy\_\allowbreak{}nj}
  attributes CPU core energy strictly to the benchmarked child PID
  (\texttt{0.073}--\texttt{0.109\ J} for native Rust/C++), Linux
  \texttt{/\allowbreak{}sys/\allowbreak{}class/\allowbreak{}powercap/\allowbreak{}intel-\allowbreak{}rapl:0} (\texttt{L2}) integrates
  socket-wide package power (\texttt{1.89}--\texttt{2.06\ J} package vs
  \texttt{0.13}--\texttt{1.83\ J} \texttt{core}), and Windows
  \texttt{NT\ 10.\allowbreak{}0.\allowbreak{}26300} blocks unprivileged user-mode Intel
  \texttt{RAPL} \texttt{MSR} reads without a ring-0 kernel driver while
  exposing exact per-PID \texttt{L3} cycle counts
  (\texttt{QueryProcessCycleTime}).
\end{itemize}

\subsection{\texorpdfstring{6.2 \texttt{GE-H2}: Three-OS and cross-ISA
25,000-tick floating-point drift (\texttt{IEEE\ 754}
\texttt{f64})}{6.2 GE-H2: Three-OS and cross-ISA 25,000-tick floating-point drift (IEEE 754 f64)}}\label{ge-h2-three-os-and-cross-isa-25000-tick-floating-point-drift-ieee-754-f64}

To test whether \texttt{IEEE\ 754} \texttt{binary64} (\texttt{f64})
simulations stay synchronized across \texttt{Windows\ 11},
\texttt{macOS}, and \texttt{Linux} (\texttt{x86\_64} and
\texttt{arm64}/\allowbreak{}\texttt{aarch64}) when \texttt{-\allowbreak{}ffast-\allowbreak{}math} and
\texttt{/fp:fast} are \emph{forbidden}, the identical 16-body chaotic
arena was stepped for \texttt{25,000} ticks
(\texttt{\textasciitilde{}6.94} minutes at \texttt{60\ Hz}), hashing
canonical \texttt{IEEE\ 754} bits (\texttt{t\_\allowbreak{}div\_\allowbreak{}ieee\_\allowbreak{}bits}),
quantized \texttt{10\^{}-3} arena-unit positions
(\texttt{t\_\allowbreak{}div\_\allowbreak{}macro\_\allowbreak{}1e3}), and discrete gameplay state
(\texttt{t\_\allowbreak{}div\_\allowbreak{}gameplay}) at every single tick:

{\scriptsize\setlength{\tabcolsep}{3.0pt}\renewcommand{\arraystretch}{1.14}
\begin{longtable}[]{@{}
  >{\raggedright\arraybackslash}p{(\linewidth - 10\tabcolsep) * \real{0.2026}}
  >{\raggedright\arraybackslash}p{(\linewidth - 10\tabcolsep) * \real{0.2222}}
  >{\raggedright\arraybackslash}p{(\linewidth - 10\tabcolsep) * \real{0.2876}}
  >{\raggedleft\arraybackslash}p{(\linewidth - 10\tabcolsep) * \real{0.0915}}
  >{\raggedleft\arraybackslash}p{(\linewidth - 10\tabcolsep) * \real{0.0915}}
  >{\raggedleft\arraybackslash}p{(\linewidth - 10\tabcolsep) * \real{0.1046}}@{}}
\toprule\noalign{}
\begin{minipage}[b]{\linewidth}\raggedright
\textbf{Host / Target A}
\end{minipage} & \begin{minipage}[b]{\linewidth}\raggedright
\textbf{Host / Target B}
\end{minipage} & \begin{minipage}[b]{\linewidth}\raggedright
\textbf{Perturbation mechanism}
\end{minipage} & \begin{minipage}[b]{\linewidth}\raggedleft
\textbf{\texttt{t\_\allowbreak{}div\_\allowbreak{}ieee\_\allowbreak{}bits}}
\end{minipage} & \begin{minipage}[b]{\linewidth}\raggedleft
\textbf{\texttt{t\_\allowbreak{}div\_\allowbreak{}macro\_\allowbreak{}1e3}}
\end{minipage} & \begin{minipage}[b]{\linewidth}\raggedleft
\textbf{\texttt{t\_\allowbreak{}div\_\allowbreak{}gameplay}}
\end{minipage} \\
\midrule\noalign{}
\endhead
\bottomrule\noalign{}
\endlastfoot
\texttt{Windows\ 11\ x86\_\allowbreak{}64} (\texttt{MSVC\ /\allowbreak{}fp:strict}, \texttt{Rust},
\texttt{V8}, \texttt{Py}) & \texttt{arm64\ macOS} \&
\texttt{x86\_64}/\allowbreak{}\texttt{aarch64\ Linux}
(\texttt{Clang}/\allowbreak{}\texttt{GCC}/\allowbreak{}\texttt{Rust}) & 3-OS \& 3-ISA \texttt{f64}
\& \texttt{f32} algebraic + \texttt{sqrt}, \texttt{FMA} off
(\texttt{/\allowbreak{}fp:strict}, \texttt{-\allowbreak{}ffp-\allowbreak{}contract=off}, no trig) &
\texttt{None\ (\textgreater{}25,000)} &
\texttt{None\ (\textgreater{}25,000)} &
\texttt{None\ (\textgreater{}25,000)} \\
\texttt{x86\_\allowbreak{}64\ Linux\ glibc\ 2.\allowbreak{}43} (\texttt{f64\_trig}) &
\texttt{aarch64\ Linux\ glibc\ 2.\allowbreak{}39} (\texttt{f64\_trig}) & Cross-ISA
Linux GNU \texttt{glibc} \texttt{std::sin/\allowbreak{}cos}
(\texttt{-\allowbreak{}ffp-\allowbreak{}contract=off}, no \texttt{-\allowbreak{}ffast-\allowbreak{}math}) &
\texttt{None\ (\textgreater{}25,000)} &
\texttt{None\ (\textgreater{}25,000)} &
\texttt{None\ (\textgreater{}25,000)} \\
\textbf{\texttt{Windows\ 11\ x86\_\allowbreak{}64\ UCRT} (\texttt{MSVC\ /\allowbreak{}fp:strict}
\& \texttt{Rust})} & \textbf{\texttt{x86\_\allowbreak{}64\ Linux\ GNU\ glibc\ 2.\allowbreak{}43}
(\texttt{Clang} \& \texttt{Rust}, \texttt{-\allowbreak{}ffp-\allowbreak{}contract=off})} &
\textbf{Cross-OS same-ISA (\texttt{x86\_64}) \texttt{std::sin/\allowbreak{}cos}
(\texttt{ucrtbase.\allowbreak{}dll} vs \texttt{glibc}, no fast-math)} &
\textbf{\texttt{10}} (\texttt{0.17\ s}) & \textbf{\texttt{7,899}} &
\textbf{\texttt{10,392}} (\texttt{173.2\ s}) \\
\textbf{Apple \texttt{arm64} macOS (\texttt{f64\_trig},
\texttt{-\allowbreak{}ffp-\allowbreak{}contract=off})} & \textbf{AMD/Intel \texttt{x86\_64} \&
Thor \texttt{aarch64} Linux (\texttt{f64\_trig},
\texttt{-\allowbreak{}ffp-\allowbreak{}contract=off})} & \textbf{Cross-OS standard \texttt{libm}
\texttt{std::sin/\allowbreak{}cos} (\texttt{Apple\ libSystem} vs \texttt{GNU\ glibc},
no \texttt{-\allowbreak{}ffast-\allowbreak{}math})} & \textbf{\texttt{10}} (\texttt{0.17\ s}) &
\textbf{\texttt{5,900}} & \textbf{\texttt{9,798}} (\texttt{163.3\ s}) \\
\textbf{\texttt{Windows\ 11\ x86\_\allowbreak{}64\ UCRT} (\texttt{MSVC\ /\allowbreak{}fp:strict}
\& \texttt{Rust})} & \textbf{Apple \texttt{arm64} macOS
(\texttt{libSystem}, \texttt{-\allowbreak{}ffp-\allowbreak{}contract=off})} & \textbf{Cross-OS \&
Cross-ISA \texttt{std::sin/\allowbreak{}cos} (\texttt{Windows\ UCRT} vs
\texttt{Apple\ libSystem}, no fast-math)} & \textbf{\texttt{37}}
(\texttt{0.62\ s}) & \textbf{\texttt{5,900}} & \textbf{\texttt{9,798}}
(\texttt{163.3\ s}) \\
\textbf{\texttt{Windows\ 11\ MSVC}
(\texttt{/\allowbreak{}O2\ /\allowbreak{}arch:AVX2\ /\allowbreak{}fp:precise} default)} &
\textbf{\texttt{x86\_\allowbreak{}64\ Linux\ Clang\ 21}
(\texttt{-\allowbreak{}O3\ -\allowbreak{}mavx2\ -\allowbreak{}mfma\ -\allowbreak{}ffp-\allowbreak{}contract=fast} default)} &
\textbf{Cross-compiler \texttt{x86\_64} \texttt{FMA} default
(\texttt{MSVC\ /\allowbreak{}fp:precise} no-contract vs \texttt{Clang} \texttt{FMA},
algebraic only)} & \textbf{\texttt{192}} (\texttt{3.20\ s}) &
\textbf{\texttt{5,927}} & \textbf{\texttt{8,869}} (\texttt{147.8\ s}) \\
\textbf{\texttt{x86\_64} Linux
(\texttt{-\allowbreak{}O3\ -\allowbreak{}mavx2\ -\allowbreak{}mfma\ -\allowbreak{}ffp-\allowbreak{}contract=fast})} & \textbf{Apple
\texttt{arm64} \& Thor \texttt{aarch64}
(\texttt{-\allowbreak{}O3\ -\allowbreak{}ffp-\allowbreak{}contract=fast})} & \textbf{Cross-ISA \texttt{FMA}
lowering (\texttt{x86\_\allowbreak{}64\ VFMADD/\allowbreak{}VFNMADD} vs
\texttt{AArch64\ FMADD/\allowbreak{}FMSUB}, algebraic only, no \texttt{-\allowbreak{}ffast-\allowbreak{}math})}
& \textbf{\texttt{406}} (\texttt{6.77\ s}) & \textbf{\texttt{7,422}} &
\textbf{\texttt{11,205}} (\texttt{186.8\ s}) \\
\textbf{\texttt{Windows\ 11} \& \texttt{Linux\ x86\_\allowbreak{}64} \texttt{Node.js}
(\texttt{V8\ Math.\allowbreak{}sin/\allowbreak{}cos})} & \textbf{\texttt{macOS\ arm64} \&
\texttt{Linux\ aarch64} \texttt{Node.js} (\texttt{V8\ Math.\allowbreak{}sin/\allowbreak{}cos})} &
\textbf{Cross-ISA Node.js V8 \texttt{Math.\allowbreak{}sin/\allowbreak{}cos} (\texttt{0} drift
Win11 vs Linux \texttt{x86\_64}; diverges vs
\texttt{arm64}/\allowbreak{}\texttt{aarch64})} & \textbf{\texttt{66}}
(\texttt{1.10\ s}) & \textbf{\texttt{7,617}} & \textbf{\texttt{10,093}}
(\texttt{168.2\ s}) \\
\textbf{\texttt{cpp\_\allowbreak{}f64\_\allowbreak{}strict} (\texttt{/\allowbreak{}fp:strict},
\texttt{-\allowbreak{}ffp-\allowbreak{}contract=off})} & \textbf{\texttt{MSVC\ /\allowbreak{}fp:fast},
\texttt{Rust\ mul\_\allowbreak{}add}, or \texttt{Clang\ -\allowbreak{}ffast-\allowbreak{}math}} &
\textbf{Same-host \texttt{FMA} contraction (\texttt{/\allowbreak{}fp:strict} vs
\texttt{/fp:fast} or explicit \texttt{mul\_add})} &
\textbf{\texttt{192}} (\texttt{3.20\ s}) &
\textbf{\texttt{5,927}--\texttt{6,576}} &
\textbf{\texttt{8,869}--\texttt{10,093}} \\
\textbf{\texttt{cpp\_\allowbreak{}f64\_\allowbreak{}strict} / \texttt{trig}} &
\textbf{\texttt{cpp\_\allowbreak{}f32\_\allowbreak{}strict} / \texttt{trig}} & \textbf{Precision
step-down (\texttt{binary64} \texttt{f64} vs \texttt{binary32}
\texttt{f32})} & \textbf{\texttt{1}} & \textbf{\texttt{8}--\texttt{9}} &
\textbf{\texttt{1,370}--\texttt{2,240}} \\
\end{longtable}
}

Four sharp 3-OS and cross-ISA findings emerge:

\begin{enumerate}
\def\labelenumi{\arabic{enumi}.}
\tightlist
\item
  \textbf{IEEE 754-2008 \texttt{+}, \texttt{-}, \texttt{*}, \texttt{/},
  and \texttt{sqrt} with \texttt{FMA} disabled (\texttt{/\allowbreak{}fp:strict} and
  \texttt{-\allowbreak{}ffp-\allowbreak{}contract=off}) are bitwise portable across
  \texttt{Windows\ 11}, \texttt{macOS}, and \texttt{Linux} on both
  \texttt{x86\_64} and \texttt{arm64}/\allowbreak{}\texttt{aarch64}} (\texttt{0}
  mismatches over \texttt{25,000} ticks across \texttt{MSVC},
  \texttt{Clang}, \texttt{GCC}, \texttt{Rustc}, \texttt{V8}, and
  \texttt{CPython}), because all five operations are mandated to round
  to \texttt{0.5} ULP.
\item
  \textbf{All three major OS math libraries (\texttt{Windows\ UCRT},
  \texttt{macOS\ libSystem}, and \texttt{Linux\ GNU\ glibc}) form a
  pairwise-divergent triangle under strict \texttt{f64}
  (\texttt{no\ fast-\allowbreak{}math})}: even on the \emph{same} \texttt{x86\_64}
  ISA, Windows \texttt{ucrtbase.\allowbreak{}dll} (\texttt{10.\allowbreak{}0.\allowbreak{}26100.\allowbreak{}9711}) diverges
  from Linux \texttt{glibc\ 2.\allowbreak{}43} at tick \texttt{10}
  (\texttt{t\_\allowbreak{}div\_\allowbreak{}gameplay\ =\ 10,\allowbreak{}392}), while \texttt{Windows\ UCRT}
  diverges from \texttt{Apple\ libSystem} at tick \texttt{37}
  (\texttt{t\_\allowbreak{}div\_\allowbreak{}gameplay\ =\ 9,\allowbreak{}798}) and \texttt{Apple\ libSystem}
  diverges from \texttt{glibc} at tick \texttt{10}
  (\texttt{t\_\allowbreak{}div\_\allowbreak{}gameplay\ =\ 9,\allowbreak{}798}). Within Windows 11,
  \texttt{MSVC\ /\allowbreak{}fp:strict}, \texttt{Rust\ x86\_\allowbreak{}64-\allowbreak{}pc-\allowbreak{}windows-\allowbreak{}msvc}, and
  \texttt{CPython\ 3.\allowbreak{}10} all link \texttt{ucrtbase.\allowbreak{}dll} and agree for
  \texttt{\textgreater{}\ 25,000} ticks (\texttt{b1fbc917\allowbreak{}88a59cff}), but
  diverge from Windows \texttt{Node.\allowbreak{}js\ V8} at tick \texttt{149}.
\item
  \textbf{Default compiler \texttt{FMA} policy diverges across
  \texttt{MSVC} and \texttt{Clang} on the same \texttt{x86\_64} ISA
  (\texttt{t\ =\ 192}), and across \texttt{x86\_64} and \texttt{AArch64}
  (\texttt{t\ =\ 406})}: \texttt{MSVC}'s default
  \texttt{/\allowbreak{}O2\ /\allowbreak{}arch:AVX2\ /\allowbreak{}fp:precise} does \emph{not} contract
  \texttt{a\ +\ b\ *\ c} (\texttt{\textgreater{}\ 25,000} ticks
  identical to \texttt{/\allowbreak{}fp:strict}), so it diverges at tick \texttt{192}
  (\texttt{t\_\allowbreak{}div\_\allowbreak{}gameplay\ =\ 8,\allowbreak{}869}) from \texttt{Clang\ 21}'s
  default \texttt{-\allowbreak{}O3\ -\allowbreak{}mavx2\ -\allowbreak{}mfma} (\texttt{-\allowbreak{}ffp-\allowbreak{}contract=fast} on by
  default). Even when both compilers enable contraction, LLVM lowers
  negated multiply-subtract differently across \texttt{x86\_64}
  (\texttt{VFMADD}/\allowbreak{}\texttt{VFNMADD}) and \texttt{AArch64}
  (\texttt{FMADD}/\allowbreak{}\texttt{FMSUB}), diverging at tick \texttt{406}
  (\texttt{t\_\allowbreak{}div\_\allowbreak{}gameplay\ =\ 11,\allowbreak{}205}).
\item
  \textbf{JavaScript V8 \texttt{Math.sin}/\allowbreak{}\texttt{Math.cos} is
  OS-independent on \texttt{x86\_64} (\texttt{Windows\ 11\ Node\ v26.\allowbreak{}13}
  matches \texttt{Linux\ x86\_\allowbreak{}64\ Node\ v22.\allowbreak{}22} for
  \texttt{\textgreater{}\ 25,000} ticks) yet ISA-dependent
  (\texttt{x86\_64} vs \texttt{arm64}/\allowbreak{}\texttt{aarch64} diverges at tick
  \texttt{66})}.
\end{enumerate}

\subsection{\texorpdfstring{6.3 \texttt{GE-H3}: Three-OS, three-API
headless \texttt{wgpu} (\texttt{DirectX\ 12} / \texttt{Metal} /
\texttt{Vulkan}) rendering
tolerance}{6.3 GE-H3: Three-OS, three-API headless wgpu (DirectX 12 / Metal / Vulkan) rendering tolerance}}\label{ge-h3-three-os-three-api-headless-wgpu-directx-12-metal-vulkan-rendering-tolerance}

To test the right-hand side of Figure 1 across all three major desktop
graphics APIs, \texttt{ge\_\allowbreak{}h3\_\allowbreak{}wgpu\_\allowbreak{}bench} (\texttt{wgpu\ v22.\allowbreak{}1.\allowbreak{}0} with
\texttt{Backends::DX12\ \textbar{}\ Backends::VULKAN\ \textbar{}\ Backends::METAL})
rendered the four canonical \texttt{Q32.32} simulation snapshots
(\texttt{ticks\ 0,\allowbreak{}\ 250,\allowbreak{}\ 1,\allowbreak{}000,\allowbreak{}\ 10,\allowbreak{}000}) into \texttt{512×512}
\texttt{RGBA8Unorm} framebuffers (\texttt{262,144} pixels,
\texttt{1,048,576} bytes/frame) via the identical two-pass WGSL shader
across \textbf{eight hardware \texttt{GPU\ ×\ API} endpoints}
(\texttt{Apple\ M3\ Ultra\ Metal},
\texttt{AMD\ Radeon\ 8060S\ Linux\ RADV\ Vulkan},
\texttt{Intel\ Iris\ Xe\ RPL-\allowbreak{}P\ Linux\ ANV\ Vulkan},
\texttt{NVIDIA\ RTX\ 4050\ Linux\ Vulkan},
\texttt{NVIDIA\ RTX\ 2060\ Win11\ DirectX\ 12},
\texttt{NVIDIA\ RTX\ 2060\ Win11\ Vulkan},
\texttt{Intel\ Iris\ Xe\ TGL\ Win11\ DirectX\ 12}, and
\texttt{Intel\ Iris\ Xe\ TGL\ Win11\ Vulkan}), alongside
\textbf{\texttt{Microsoft\ Basic\ Render\ Driver}
(\texttt{Win11\ DirectX\ 12\ WARP})}, \textbf{\texttt{Mesa\ llvmpipe}
(\texttt{Linux\ CPU\ Vulkan})}, and the zero-float
\textbf{\texttt{CPU\ Q16.\allowbreak{}16}} integer reference rasterizer:

{\scriptsize\setlength{\tabcolsep}{2.4pt}\renewcommand{\arraystretch}{1.12}
\begin{longtable}[]{@{}
  >{\raggedright\arraybackslash}p{(\linewidth - 16\tabcolsep) * \real{0.1686}}
  >{\raggedright\arraybackslash}p{(\linewidth - 16\tabcolsep) * \real{0.1802}}
  >{\raggedleft\arraybackslash}p{(\linewidth - 16\tabcolsep) * \real{0.1221}}
  >{\raggedleft\arraybackslash}p{(\linewidth - 16\tabcolsep) * \real{0.1453}}
  >{\raggedleft\arraybackslash}p{(\linewidth - 16\tabcolsep) * \real{0.0814}}
  >{\raggedleft\arraybackslash}p{(\linewidth - 16\tabcolsep) * \real{0.0872}}
  >{\raggedleft\arraybackslash}p{(\linewidth - 16\tabcolsep) * \real{0.0814}}
  >{\raggedleft\arraybackslash}p{(\linewidth - 16\tabcolsep) * \real{0.0640}}
  >{\centering\arraybackslash}p{(\linewidth - 16\tabcolsep) * \real{0.0698}}@{}}
\toprule\noalign{}
\begin{minipage}[b]{\linewidth}\raggedright
\textbf{Adapter A (\texttt{OS} · \texttt{API} · Driver)}
\end{minipage} & \begin{minipage}[b]{\linewidth}\raggedright
\textbf{Adapter B (\texttt{OS} · \texttt{API} · Driver)}
\end{minipage} & \begin{minipage}[b]{\linewidth}\raggedleft
\textbf{Mean differing pixels (\texttt{\%})}
\end{minipage} & \begin{minipage}[b]{\linewidth}\raggedleft
\textbf{Max \(\lvert \Delta \rvert\) (\texttt{0..255})}
\end{minipage} & \begin{minipage}[b]{\linewidth}\raggedleft
\textbf{Mean MAE (\texttt{u8})}
\end{minipage} & \begin{minipage}[b]{\linewidth}\raggedleft
\textbf{Mean RMSE (\texttt{u8})}
\end{minipage} & \begin{minipage}[b]{\linewidth}\raggedleft
\textbf{Min \texttt{PSNR} (\texttt{dB})}
\end{minipage} & \begin{minipage}[b]{\linewidth}\raggedleft
\textbf{Min \texttt{SSIM}}
\end{minipage} & \begin{minipage}[b]{\linewidth}\centering
\textbf{\texttt{GE-H3} Gate}
\end{minipage} \\
\midrule\noalign{}
\endhead
\bottomrule\noalign{}
\endlastfoot
\textbf{NVIDIA RTX 2060 (\texttt{Win11\ DirectX\ 12} · \texttt{591.74})}
& \textbf{NVIDIA RTX 4050 (\texttt{Linux\ Vulkan} · \texttt{595.99.02})}
& \textbf{\texttt{0.000\%}} (\texttt{0} px) & \textbf{\texttt{0}} &
\texttt{0.00000} & \texttt{0.00000} & \textbf{\texttt{∞} (\texttt{4/4}
exact)} & \textbf{\texttt{1.000000}} & \textbf{Pass} \\
\textbf{NVIDIA RTX 2060 (\texttt{Win11\ DirectX\ 12} ·
\texttt{HLSL}/\allowbreak{}\texttt{DXIL})} & \textbf{NVIDIA RTX 2060
(\texttt{Win11\ Vulkan} · \texttt{SPIR-V})} \emph{(same GPU)} &
\textbf{\texttt{0.0002\%}} (\texttt{0–1} px) & \textbf{\texttt{1}} &
\texttt{0.000001} & \texttt{0.00049} & \textbf{\texttt{108.34}} &
\textbf{\texttt{1.000000}} & \textbf{Pass} \\
\textbf{Intel Iris Xe (\texttt{Win11\ DirectX\ 12} ·
\texttt{31.\allowbreak{}0.\allowbreak{}101.\allowbreak{}3222})} & \textbf{Intel Iris Xe (\texttt{Linux\ Vulkan}
· Mesa \texttt{ANV} 26.0.8)} & \texttt{0.0003\%} & \texttt{1} &
\texttt{0.000001} & \texttt{0.00059} & \textbf{\texttt{105.33}} &
\textbf{\texttt{1.000000}} & \textbf{Pass} \\
\textbf{Intel Iris Xe (\texttt{Win11\ DirectX\ 12} ·
\texttt{HLSL}/\allowbreak{}\texttt{DXIL})} & \textbf{Intel Iris Xe
(\texttt{Win11\ Vulkan} · \texttt{SPIR-V})} \emph{(same GPU)} &
\textbf{\texttt{0.283\%}} (\texttt{729–758} px) & \textbf{\texttt{1}} &
\texttt{0.00071} & \texttt{0.02665} & \textbf{\texttt{79.52}} &
\textbf{\texttt{0.999996}} & \textbf{Pass} \\
\textbf{Apple M3 Ultra (\texttt{macOS\ Metal} 4)} & \textbf{Intel Iris
Xe (\texttt{Win11\ DirectX\ 12} / \texttt{Linux\ ANV\ Vulkan})} &
\texttt{1.301\%} & \texttt{2} & \texttt{0.00386} & \texttt{0.06215} &
\textbf{\texttt{72.06}} & \textbf{\texttt{0.999986}} & \textbf{Pass} \\
\textbf{Apple M3 Ultra (\texttt{macOS\ Metal} 4)} & \textbf{Intel Iris
Xe (\texttt{Win11\ Vulkan} · \texttt{31.\allowbreak{}0.\allowbreak{}101.\allowbreak{}3222})} & \texttt{1.573\%}
& \texttt{2} & \texttt{0.00455} & \texttt{0.06743} &
\textbf{\texttt{71.39}} & \textbf{\texttt{0.999982}} & \textbf{Pass} \\
\textbf{AMD Radeon 8060S (\texttt{Linux\ Vulkan} · Mesa \texttt{RADV})}
& \textbf{Apple M3 Ultra (\texttt{macOS\ Metal} 4)} & \texttt{3.103\%} &
\texttt{1} & \texttt{0.00788} & \texttt{0.08876} &
\textbf{\texttt{69.09}} & \textbf{\texttt{0.999956}} & \textbf{Pass} \\
\textbf{AMD Radeon 8060S (\texttt{Linux\ Vulkan} · Mesa \texttt{RADV})}
& \textbf{Intel Iris Xe (\texttt{Win11\ DirectX\ 12} /
\texttt{Win11\ Vulkan})} & \texttt{3.774\%} / \texttt{3.565\%} &
\texttt{2} & \texttt{0.01022} / \texttt{0.00968} & \texttt{0.10116} /
\texttt{0.09844} & \textbf{\texttt{67.97}} / \textbf{\texttt{68.20}} &
\textbf{\texttt{0.999951}} & \textbf{Pass} \\
\textbf{AMD Radeon 8060S (\texttt{Linux\ Vulkan} · Mesa \texttt{RADV})}
& \textbf{NVIDIA RTX 2060 (\texttt{Win11\ DX12} / \texttt{Vulkan}) \&
RTX 4050} & \texttt{6.554\%} & \texttt{2} & \texttt{0.01735} &
\texttt{0.13180} & \textbf{\texttt{65.64}} & \textbf{\texttt{0.999874}}
& \textbf{Pass} \\
\textbf{NVIDIA RTX 2060 (\texttt{Win11\ DirectX\ 12})} & \textbf{Intel
Iris Xe (\texttt{Win11\ DirectX\ 12})} \emph{(same Win11 host)} &
\texttt{8.279\%} & \texttt{2} & \texttt{0.02149} & \texttt{0.14667} &
\textbf{\texttt{64.79}} & \textbf{\texttt{0.999844}} & \textbf{Pass} \\
\textbf{Apple M3 Ultra (\texttt{macOS\ Metal} 4)} & \textbf{NVIDIA RTX
2060 (\texttt{Win11\ DX12} / \texttt{Vulkan}) \& RTX 4050} &
\texttt{9.337\%} & \texttt{2} & \texttt{0.02467} & \texttt{0.15721} &
\textbf{\texttt{64.16}} & \textbf{\texttt{0.999833}} & \textbf{Pass} \\
\textbf{Apple M3 Ultra (\texttt{macOS\ Metal} 4)} &
\textbf{\texttt{Microsoft\ WARP} (\texttt{Win11\ DirectX\ 12} CPU)} &
\texttt{0.003\%} & \texttt{1} & \texttt{0.000008} & \texttt{0.00271} &
\textbf{\texttt{98.79}} & \textbf{\texttt{1.000000}} & \textbf{Pass} \\
\textbf{\texttt{Microsoft\ WARP} (\texttt{Win11\ DirectX\ 12} CPU)} &
\textbf{\texttt{Mesa\ llvmpipe} (\texttt{Linux\ LLVM\ 21.\allowbreak{}1.\allowbreak{}8} CPU
\texttt{Vulkan})} & \texttt{1.591\%} & \texttt{2} & \texttt{0.00543} &
\texttt{0.07359} & \textbf{\texttt{70.35}} & \textbf{\texttt{0.999983}}
& \textbf{Pass} \\
Hardware GPUs (\texttt{8} endpoints across
\texttt{DX12}/\allowbreak{}\texttt{Metal}/\allowbreak{}\texttt{Vulkan}) &
\textbf{\texttt{CPU\ Q16.\allowbreak{}16} Zero-Float Integer Reference (\texttt{4}
hosts)} & \texttt{20.400\%} -- \texttt{27.626\%} & \texttt{2} &
\texttt{0.05612} -- \texttt{0.07864} & \texttt{0.23689} --
\texttt{0.28068} & \texttt{59.13} -- \texttt{60.58} & \texttt{0.999556}
-- \texttt{0.999693} & \textbf{Pass} \\
\end{longtable}
}

\begin{itemize}
\tightlist
\item
  \textbf{Isolating shader-compiler (\texttt{DX12}
  \texttt{HLSL}/\allowbreak{}\texttt{DXIL} vs \texttt{Vulkan} \texttt{SPIR-V}) drift
  from GPU silicon}: Comparing the \textbf{exact same physical GPU on
  the exact same Windows 11 machine} across \texttt{DirectX\ 12} and
  \texttt{Vulkan} separates compiler translation (\texttt{Naga}
  \texttt{WGSL\ -\textgreater{}\ HLSL} vs
  \texttt{WGSL\ -\textgreater{}\ SPIR-V}) from hardware SFU/ROP design.
  On \texttt{NVIDIA\ GeForce\ RTX\ 2060}, \texttt{Win11\ DirectX\ 12}
  and \texttt{Win11\ Vulkan} produce bit-identical \texttt{SHA-256}
  framebuffers at ticks \texttt{0} and \texttt{250}, then diverge by
  \texttt{1} LSB on a single grazing specular pixel at ticks
  \texttt{1,000} and \texttt{10,000} (\texttt{0.0002\%} mean diff,
  \texttt{108.\allowbreak{}34\ dB} \texttt{PSNR}) --- while
  \texttt{RTX\ 2060\ (\allowbreak{}Win11\ DirectX\ 12)} matches
  \texttt{RTX\ 4050\ (\allowbreak{}Linux\ Vulkan)} \texttt{4\ /\ 4} frames
  bit-for-bit (\texttt{0} differing pixels across Turing vs Ada Lovelace
  and Windows \texttt{DX12} vs Linux \texttt{Vulkan})! On
  \texttt{Intel\ Iris\ Xe\ Graphics} (\texttt{Tiger\ Lake}),
  \texttt{Win11\ DirectX\ 12} vs \texttt{Win11\ Vulkan} differs on
  \texttt{0.283\%} of pixels (\texttt{79.52\ dB} \texttt{PSNR},
  \(\Delta_{\max} = 1\)) across all \texttt{4} frames, whereas
  \texttt{Win11\ DirectX\ 12} vs \texttt{Linux\ Mesa\ ANV\ Vulkan}
  differs on just \texttt{0.0003\%} (\texttt{105.\allowbreak{}33\ dB}, bit-identical
  on ticks \texttt{0} and \texttt{1,000}).
\item
  \textbf{Bounded tolerance confirmed across all \texttt{28\ /\ 28}
  hardware \texttt{GPU\ ×\ API} pairs (\texttt{112\ /\allowbreak{}\ 112} frames)}:
  \texttt{27\ /\ 28} hardware endpoint pairs differ in raw
  \texttt{RGBA8} \texttt{SHA-256} digests, yet \textbf{zero pixels on
  any hardware pair across \texttt{DirectX\ 12}, \texttt{Metal}, and
  \texttt{Vulkan} differed by more than \texttt{2\ /\ 255}}
  (\(p(|\Delta| > 8) = 0.000000 \le 0.001\)), \texttt{PSNR} exceeded
  \texttt{64.16\ dB} (\texttt{\textgreater{}\ 40.0\ dB} gate by
  \texttt{+24.\allowbreak{}16\ dB}), and \texttt{SSIM} exceeded \texttt{0.999833}
  (\texttt{\textgreater{}\ 0.999} gate) at
  \texttt{0.29}--\texttt{6.04\ ms} per frame.
\item
  \textbf{Software rasterizer contrast (\texttt{WARP} vs
  \texttt{llvmpipe} vs \texttt{CPU\ Q16.\allowbreak{}16})}: Even software CPU GPU-API
  emulators diverge across APIs ---
  \texttt{Microsoft\ Basic\ Render\ Driver}
  (\texttt{Win11\ DirectX\ 12\ WARP}) vs \texttt{Mesa\ llvmpipe}
  (\texttt{Linux\ Vulkan}) differs on \texttt{1.591\%} of pixels
  (\texttt{70.35\ dB} \texttt{PSNR}, \(\Delta_{\max} = 2\)) --- whereas
  the zero-float \texttt{CPU\ Q16.\allowbreak{}16} integer software rasterizer
  achieved \textbf{\texttt{0} \texttt{SHA-256} framebuffer mismatches
  across all four OS/ISA hosts} (\texttt{macOS\ arm64},
  \texttt{Linux\ AMD\ x86\_\allowbreak{}64}, \texttt{Linux\ Intel\ x86\_\allowbreak{}64}, and
  \texttt{Windows\ 11\ x86\_\allowbreak{}64}: \texttt{9c693a38.\allowbreak{}.\allowbreak{}.\allowbreak{}} at tick
  \texttt{0}, \texttt{ce3b2cec.\allowbreak{}.\allowbreak{}.\allowbreak{}} at tick \texttt{10,000}).
\end{itemize}

\subsection{\texorpdfstring{6.4 Playable joule-native \texttt{Q32.32}
Chrono-Budget Arena
(\texttt{/\allowbreak{}living/\allowbreak{}game-\allowbreak{}engines/\allowbreak{}})}{6.4 Playable joule-native Q32.32 Chrono-Budget Arena (/living/game-engines/)}}\label{playable-joule-native-q32.32-chrono-budget-arena-livinggame-engines}

Joining this study's exact \texttt{Q32.32} replay contract with the
energy-native temporal mechanics of
\href{https://research.openie.dev/papers/time-as-an-index/}{Time as an
Index} (Section 9.2), Instrument 4 of
\href{https://research.openie.dev/living/game-engines/}{\texttt{/\allowbreak{}living/\allowbreak{}game-\allowbreak{}engines/\allowbreak{}}}
ships a playable \texttt{60\ Hz} deterministic orbital-navigation and
pickup-defense mini-game running in JavaScript \texttt{BigInt}
(\texttt{0} floating-point operations inside simulation):

\begin{itemize}
\tightlist
\item
  \textbf{Live \texttt{SHA-256} state tape \& deterministic replay}:
  Every tick updates a rolling \texttt{SHA-256} state digest and records
  compact 1-byte player input events into \(D\); clicking
  \texttt{Replay\ \&\ Verify\ SHA-\allowbreak{}256} re-executes the entire match from
  \texttt{(\allowbreak{}seed,\allowbreak{}\ D)} and verifies bitwise hash equality
  (\texttt{replayed\_\allowbreak{}sha256\ ===\ live\_\allowbreak{}sha256}).
\item
  \textbf{Joule-budgeted temporal lookahead \& rewind (\texttt{est\_j}
  design model)}: Querying future trajectories
  (\texttt{Hellman\ Branch\ Preview}, simulating \texttt{3} candidate
  steering branches \texttt{K} ticks ahead) or rewinding to a persistent
  causal checkpoint debits an explicit microjoule (\texttt{µJ}) battery
  (\texttt{step\_\allowbreak{}cost\_\allowbreak{}uj} per speculative tick plus Landauer
  branch-erasure cost when non-chosen futures or rewound timelines are
  discarded), playable via keyboard, mouse/touch, and \texttt{WebXR} /
  \texttt{Gamepad} controllers (including Meta Quest Browser).
\end{itemize}

\section{7. Infeasibility flags}\label{infeasibility-flags}

\begin{enumerate}
\def\labelenumi{\arabic{enumi}.}
\tightlist
\item
  \textbf{Cross-GPU and cross-API bit-exact hardware rendering:
  infeasible (\texttt{27\ /\ 28} hardware \texttt{GPU\ ×\ API} pairs ---
  including \texttt{2\ /\ 2} same-GPU \texttt{DirectX\ 12} vs
  \texttt{Vulkan} pairs --- diverge in raw \texttt{SHA-256} in
  \texttt{GE-H3}), whereas bounded perceptual tolerance
  (\texttt{PSNR\ \textgreater{}=\ 64.16\ dB},
  \texttt{SSIM\ \textgreater{}=\ 0.99983}, \(\Delta_{\max} \le 2\))
  holds across \texttt{DirectX\ 12}, \texttt{Metal}, and \texttt{Vulkan}
  when GPU outputs never feed back into \(S\).}
\item
  \textbf{Cross-platform exact physics with stock floating-point
  settings: infeasible across \texttt{Windows}, \texttt{macOS}, and
  \texttt{Linux}.} Even \emph{without} \texttt{-\allowbreak{}ffast-\allowbreak{}math} or
  \texttt{/fp:fast}, standard \texttt{std::sin}/\allowbreak{}\texttt{std::cos} forms
  a 3-way divergent triangle across \texttt{Windows\ UCRT},
  \texttt{Apple\ libSystem}, and \texttt{GNU\ glibc}
  (\texttt{t\_\allowbreak{}div\ =\ 10} to \texttt{37}), default
  \texttt{MSVC\ /\allowbreak{}fp:precise} vs \texttt{Clang\ -\allowbreak{}ffp-\allowbreak{}contract=fast}
  diverges on \texttt{x86\_64} at \texttt{t\_\allowbreak{}div\ =\ 192}, and cross-ISA
  \texttt{FMA} (\texttt{x86\_64} vs \texttt{arm64}) diverges at
  \texttt{t\_\allowbreak{}div\ =\ 406}, cascading into gameplay desynchronization
  within \texttt{8,869}--\texttt{11,205} ticks (\texttt{GE-H2}).
\item
  \textbf{GPU compute inside the simulation: infeasible as exact;
  feasible only with deterministic reductions} at a measured cost of
  about 0.2\% to 7.8\% on the cited hardware, and still subject to
  compiler contraction unless guarded per instruction (SPIR-V
  NoContraction) on every target.
\item
  \textbf{Engine-portable behavior: not demonstrated in any external
  source read; demonstrated at the 2D physics kernel level in
  \texttt{GE-H1} across \texttt{30} target × host cells (\texttt{6}
  machines, \texttt{3} OSs
  \texttt{macOS}/\allowbreak{}\texttt{Linux}/\allowbreak{}\texttt{Windows\ 11}, \texttt{3} ISAs,
  \texttt{7} runtimes/compilers), while full engine projections
  (\texttt{GE-H4}, \texttt{GE-H9}) remain open.}
\item
  \textbf{Interpreted rules at large actor counts: constrained.} The
  GameRule measurement implies a ceiling near 300 actors at 20 ms
  (\texttt{GE-H1} further shows Python \texttt{int} costs
  \texttt{73×}--\texttt{132×} native Rust \texttt{reported\_\allowbreak{}j} on
  Linux/macOS and \texttt{179×} wall time on Windows 11); a compiled
  native or WebAssembly projection is required at scale.
\item
  \textbf{Proprietary third-party game engines as exact targets:
  unverified.} While bare-metal Windows 11 \texttt{PE} (\texttt{rustc}
  and \texttt{MSVC}) and \texttt{DirectX\ 12} (\texttt{wgpu}) execution
  are verified in Section 6, Unity's manual, as quoted by GameRule,
  disclaims ordering guarantees; the Unity and Unreal template catalogs
  could not be read (their pages render client-side), and no Unreal
  execution-order source was read.
\end{enumerate}

\section{8. Sources and limits}\label{sources-and-limits}

Every source above was read in full at the passage cited, except: the
GameCraft-Bench paper (abstract and repository README read), the RCE
about page (summary page read), and the Capcom press story (one
article). Companion house studies cited for telemetry levels
(\texttt{SATB}) and joule-budgeted temporal game mechanics are
\href{https://research.openie.dev/papers/time-as-an-index/}{Time as an
Index} and the
\href{https://research.openie.dev/papers/satb-harness/}{Satiation Tiered
Benchmark Harness} (\texttt{{[}same\ house{]}}). Sources consulted but
not cited as evidence: the \href{https://ludii.games/}{Ludii portal},
which describes games as structured sets of ludemes; the Game Grammar
hobby repository; and the
\href{https://github.com/playcanvas/engine}{PlayCanvas},
\href{https://github.com/BabylonJS/Babylon.js}{Babylon.js} and
\href{https://github.com/mrdoob/three.js}{three.js} READMEs, which
describe rendering libraries rather than game description formats.

Not read, paywalled or otherwise unavailable: T. Schaul, ``A Video Game
Description Language for Model-based or Interactive Learning,'' IEEE CIG
2013,
\href{https://doi.org/10.1109/CIG.2013.6633610}{doi:10.1109/CIG.2013.6633610};
J. Gregory, \emph{Game Engine Architecture},
\href{https://doi.org/10.1201/9781315267845}{doi:10.1201/9781315267845}.
M. Ebner et al., ``Towards a Video Game Description Language,''
\href{https://doi.org/10.4230/DFU.Vol6.12191.85}{doi:10.4230/DFU.Vol6.12191.85},
is open access but was not read for this study.

In Section 6, \texttt{GE-H1}, \texttt{GE-H2}, and \texttt{GE-H3} leave
\texttt{measured\_\allowbreak{}j\ =\ None} and \texttt{board\_\allowbreak{}metered\ =\ false}
across all six machines (separating macOS \texttt{L3} per-process
\texttt{ri\_\allowbreak{}energy\_\allowbreak{}nj}, Linux \texttt{L2} socket/core \texttt{RAPL}
\texttt{reported\_\allowbreak{}j}, and Windows 11 \texttt{L3}
\texttt{GetProcessTimes}/\allowbreak{}\texttt{QueryProcessCycleTime} activity
counters, and keeping the Chrono-Budget Arena's \texttt{est\_j}
mechanics strictly analytical). Full multi-engine scene/behavior
emitters (\texttt{GE-H4}--\texttt{GE-H9}) remain open work.


\end{document}
